% Options for packages loaded elsewhere
\PassOptionsToPackage{unicode}{hyperref}
\PassOptionsToPackage{hyphens}{url}
\documentclass[
]{article}
\usepackage{xcolor}
\usepackage{amsmath,amssymb}
\setcounter{secnumdepth}{-\maxdimen} % remove section numbering
\usepackage{iftex}
\ifPDFTeX
  \usepackage[T1]{fontenc}
  \usepackage[utf8]{inputenc}
  \usepackage{textcomp} % provide euro and other symbols
\else % if luatex or xetex
  \usepackage{unicode-math} % this also loads fontspec
  \defaultfontfeatures{Scale=MatchLowercase}
  \defaultfontfeatures[\rmfamily]{Ligatures=TeX,Scale=1}
\fi
\usepackage{lmodern}
\ifPDFTeX\else
  % xetex/luatex font selection
\fi
% Use upquote if available, for straight quotes in verbatim environments
\IfFileExists{upquote.sty}{\usepackage{upquote}}{}
\IfFileExists{microtype.sty}{% use microtype if available
  \usepackage[]{microtype}
  \UseMicrotypeSet[protrusion]{basicmath} % disable protrusion for tt fonts
}{}
\makeatletter
\@ifundefined{KOMAClassName}{% if non-KOMA class
  \IfFileExists{parskip.sty}{%
    \usepackage{parskip}
  }{% else
    \setlength{\parindent}{0pt}
    \setlength{\parskip}{6pt plus 2pt minus 1pt}}
}{% if KOMA class
  \KOMAoptions{parskip=half}}
\makeatother
\usepackage{longtable,booktabs,array}
\newcounter{none} % for unnumbered tables
\usepackage{calc} % for calculating minipage widths
% Correct order of tables after \paragraph or \subparagraph
\usepackage{etoolbox}
\makeatletter
\patchcmd\longtable{\par}{\if@noskipsec\mbox{}\fi\par}{}{}
\makeatother
% Allow footnotes in longtable head/foot
\IfFileExists{footnotehyper.sty}{\usepackage{footnotehyper}}{\usepackage{footnote}}
\makesavenoteenv{longtable}
\setlength{\emergencystretch}{3em} % prevent overfull lines
\providecommand{\tightlist}{%
  \setlength{\itemsep}{0pt}\setlength{\parskip}{0pt}}
\usepackage{bookmark}
\IfFileExists{xurl.sty}{\usepackage{xurl}}{} % add URL line breaks if available
\urlstyle{same}
\hypersetup{
  hidelinks,
  pdfcreator={LaTeX via pandoc}}

\author{}
\date{}

\begin{document}

\section{\texorpdfstring{\textbf{The Geometry of
Failure}}{The Geometry of Failure}}\label{the-geometry-of-failure}

\textbf{Taylor William Buley}

Independent Researcher

https://buley.fyi/

\subsection{\texorpdfstring{\textbf{Abstract}}{Abstract}}\label{abstract}

I model \textbf{fork/race/fold} as a reusable computational primitive:
fork work into parallel streams, race streams to select valid progress,
fold results through deterministic reconciliation, and vent paths whose
continued existence would destabilize the whole. The central claim of
this manuscript is that failure in such systems is not a binary event
but a \textbf{topological coordinate}.

Failure is not the opposite of success; in the geometry of distributed
systems, it is its mathematical guarantor. Given the broad scope of
distributed systems -- everything from cellular replication to global
financial markets -- the impact of the subsequent proof of this truth is
sweeping.

I present a geometric description of \textbf{stability}. Fork introduces
curvature into the computation graph by raising the first Betti number
!{[}{]}{[}image1{]}. Race moves the state along that curved manifold.
Fold projects the manifold back toward a stable sink. Vent does not
merely discard work; it maps the \textbf{Void} -- the bounded region of
lost branches, unpaid debt, and discarded alternatives that every finite
collapse must account for.

I report structural similarities in selected natural and engineered
examples: \emph{Physarum polycephalum} recreated a rail-like network
over nutrient gradients

!{[}{]}{[}image2{]}, myelinated neurons pipeline action potentials (with
measured large speedups), photosynthetic antenna complexes exhibit high
step-level exciton-transfer efficiency in cited systems, and DNA
replication uses out-of-order fragment synthesis with deterministic
reassembly (Okazaki fragments). In each case, the interesting question
is not simply how much work is dissipated, but whether the system's
transition kernel points back toward a bounded stable region after
perturbation.

This manuscript's contribution is treating these systems as one
operational abstraction with explicit diagnostics and executable
obligations. I present the \textbf{Wallington Rotation}, a scheduling
algorithm that rotates partially ordered work into concurrent
stage-local tracks with controlled reconciliation, and I show
(constructively plus executable verification) that four primitives --
fork, race, fold, vent -- are sufficient for the finite DAG classes used
in this paper's implementation scope. I give the algorithm a natural
\textbf{topological characterization}: fork increases
!{[}{]}{[}image1{]} (creating independent parallel paths), race
traverses homotopy-equivalent paths simultaneously, fold projects
!{[}{]}{[}image1{]} back toward zero, and vent propagation is a natural
transformation that releases unstable descendants while preserving
sibling structure. Self-describing frames create a \textbf{covering
space} over the computation graph -- working in the cover (multiplexed,
out-of-order) then projecting back to the base space (sequential,
reassembled).

The central analytic objects are \textbf{curvature} and \textbf{drift} .
A Foster-Lyapunov drift schema asks whether, for states outside a
bounded small set, the expected motion of the system points inward. In
this framing, failure is the moment that restorative pressure is lost:
arrival pressure exceeds the combined recovery pressure , the drift
field turns outward, and the failure manifold expands. Recovery is the
negative drift back toward the stable sink.

I then show that \textbf{selected canonical queueing constructions are
recovered at the} \textbf{boundary in this framework}. Little's Law and
Jackson-style queueing results are treated as path-like boundary
examples in the modeled scope

. The \textbf{pipeline Reynolds number} extends the analysis with
topology-sensitive regimes (laminar, transitional, turbulent) and
doubles as a gauge of failure turbulence. I also use quantum-mechanical
terminology -- superposition, tunneling, interference, entanglement,
measurement, collapse -- as structural correspondence language for
mapped computational operations. I derive the \textbf{topological
deficit} as a stability diagnostic: in the analyzed examples, matched
Betti structure () correlates with higher fit/efficiency, while
co-occurs with measurable waste and brittleness (healthcare delays,
settlement lockup, protocol-level blocking)

. I define the \textbf{Bule} (1 B = 1 unit of ) as the corresponding
unit of topological stress.

These cross-domain correspondences are exemplar-based and correlational;
they are not presented as a systematic causal survey or universal proof.
In compression benchmarks on homogeneous web content, standalone global
brotli retained better ratio, so the topological claim is strategy
subsumption, framing reduction, and portability rather than universal
ratio superiority. The stronger conclusion is geometric:
\textbf{reliability is geometry}, failure is the void left by unmanaged
topological deficit, and validation should prove negative drift rather
than merely count broken cases.

o -\textgreater{} o -\textgreater{} o -\textgreater{} o -\textgreater{}
o -\textgreater{} o

The conveyor belt was not new when Ford adopted it in 1913. It is a path
graph -- a line: one-dimensional, simply connected and without
branching, where interior nodes have one predecessor and one successor.
Modern pipelines can optimize this structure, but in this paper it is
treated as a boundary case of a richer topology class.

\begin{verbatim}
  o \-\> o  
 /     \\  
\end{verbatim}

o -\textgreater{} o o -\textgreater{} o\\
\textbackslash{} /\\
o -\textgreater{} o

Fork/race/fold is represented here as a directed acyclic graph (DAG)
with merge points: nodes branch, paths run in parallel and merge
vertices fold concurrent paths into one. In the analyzed domains,
recurring bottlenecks arise when high- workloads are forced through
path-like structures. The operational remedy is to work in the cover
space (multiplexed, out-of-order) and project back to the base space
(sequential, reassembled).

I instantiate the algorithm in \textbf{five} domains, presented in stack
order --- from building blocks to bytes on wire --- each layer enabled
by the ones below it.

\begin{enumerate}
\def\labelenumi{\arabic{enumi}.}
\tightlist
\item
  In formal verification (the foundation, §10), I implement a temporal
  logic model checker (@affectively/aeon-logic) whose BFS state-space
  exploration is itself a fork/race/fold computation. Each
  multi-successor expansion is a fork, each transition to an
  already-visited state is a fold (interference), each unfair cycle
  filtered by weak fairness is a vent, and termination is collapse. I
  wrap these invariants in a \textbf{Failure-Modulated Flow Validation
  (FMFV)} schema: not only must and hold, but the expected drift must
  point back toward a bounded small set. Stability becomes a proof of
  inward pressure rather than a count of non-violations\\
  .\\
\item
  In formal language theory (the programming model, §11), I implement
  Gnosis\\
  --- a programming language whose source code \emph{is} the computation
  graph and whose compiler \emph{is} a fork/race/fold pipeline. Programs
  are Cypher-like graphs with four edge types (FORK, RACE, FOLD, VENT),
  carrying \textbf{drift coefficients} that declare restorative or
  destabilizing pressure on each edge. The compiler goes beyond counting
  ; Betti verifies the routing operator's \textbf{spectral radius} and
  rejects programs whose local flow cannot return to a stable sink.\\
\item
  In distributed staged computation (the scheduling algorithm, §7),
  chunked pipelined processing reduces sequential depth from to ,
  yielding modeled step-count speedups of 3.1x--267x in the listed
  scenarios under explicit idealized assumptions (uniform stage service
  time, zero inter-node communication cost). The Wallington Rotation
  functions as not only a speedup law but a \textbf{stability law}: it
  reshapes the manifold so flow keeps recovering after local drops. The
  Worthington Whip acts as the controlled payment of reconciliation
  tension.\\
\item
  In edge transport (the wire format, §8), I implement a binary stream
  protocol with 10-byte self-describing frame headers and native
  fork/race/fold operations on UDP, reducing framing overhead by 95
  percent versus HTTP/1.1 and removing one-path ordered-delivery
  coupling that drives head-of-line behavior in TCP-bound stacks. The
  header becomes the \textbf{coordinate chart} of the cover space.
  Head-of-line blocking is interpreted as a topological contradiction:
  independent paths forced to share a single dimension.\\
\item
  In compression (bytes on wire, §9), I implement per-chunk topological
  codec racing (fork codecs, race per chunk, fold to winner), with
  executable verification of roundtrip correctness, codec-vent behavior
  and invariants\\
  . The capstone is not merely finding the best ratio; it is exposing
  the \textbf{Pareto frontier} between ratio, portability, random
  access, and failure cost. The two-level race searches for the minimal
  cost floor at which the system can collapse to one encoding without
  losing control of the manifold.
\end{enumerate}

Within the modeled scope in this paper (finite DAG decompositions under
C1-C4), the algorithm is a high-fit topology class with measurable fit
via . It is intentionally simple: four primitives, explicit assumptions,
and executable checks.

The technique and tooling provide a method to identify, measure and
reduce topological waste in high-impact domains, including drug
discovery, health care and energy systems.

\subsection{\texorpdfstring{\textbf{0. A Child, a Ball, a
Line}}{0. A Child, a Ball, a Line}}\label{a-child-a-ball-a-line}

Imagine a child handing a ball to a friend in a line. Four children, one
hundred balls. The first child hands Ball 1 to the second, waits while
it travels through all four kids, then hands Ball 2. Everyone stands
idle while one ball moves. Four hundred handoffs, one at a time.

Now imagine something slightly different. The moment the first child
passes Ball 1 to the second, she picks up Ball 2. Now the second child
passes Ball 1 to the third while the first child passes Ball 2 to the
second. Everyone is busy at once. One hundred balls, four children, one
hundred and \textbf{three} handoffs. This is \textbf{pipelining} -- a
known technique in computer architecture that has been used since the
1960s.

But what if the first child could juggle? Bundle twenty-five balls
together, pass them as a single armful. One hundred balls, four
children, chunks of twenty-five: \textbf{seven handoffs}. That is a 98
percent reduction. This is \textbf{chunked pipelining}, and the formula
is:

where is the number of balls, is the chunk size and is the number of
children.

Handoffs are unnecessary overhead: waste to be eliminated. In real life,
they take the form of network packets, memory allocations, context
switches, and other system resources that consume time and energy. They
are also computationally ugly, as we show in section 2.

\subsubsection{\texorpdfstring{\textbf{0.1 The
Triangle}}{0.1 The Triangle}}\label{the-triangle}

Now look at what happens when four chunks move through four children.
Draw it as a grid where time moves left to right, children are stacked
top to bottom. This is the \textbf{grid}:

Time: t1 t2 t3 t4 t5 t6 t7\\
Kid 1: {[}C1{]} {[}C2{]} {[}C3{]} {[}C4{]}\\
Kid 2: {[}C1{]} {[}C2{]} {[}C3{]} {[}C4{]}\\
Kid 3: {[}C1{]} {[}C2{]} {[}C3{]} {[}C4{]}\\
Kid 4: {[}C1{]} {[}C2{]} {[}C3{]} {[}C4{]}

Focus on the \textbf{ramp-up} -- the first four time steps. Read what's
active at each moment. This is the \textbf{triangle}:

t1: 1\\
t2: 2 1\\
t3: 3 2 1\\
t4: 4 3 2 1

A triangle. The top has one chunk. Each row adds one more. At the
pipeline is full -- every child is busy. Now trace any path through this
triangle:

\begin{itemize}
\tightlist
\item
  \textbf{Read any column} (one child's work over time): 1, 2, 3, 4.
  Correct order.\\
\item
  \textbf{Read any row} (all children at one moment): a contiguous
  subsequence, in order under the stated stage/sequence dependencies.\\
\item
  \textbf{Read the diagonal} (, all four children active): 4, 3, 2, 1 --
  the wavefront. Every chunk is at a different stage, but they are all
  progressing in the correct relative order.
\end{itemize}

\textbf{Under this dependency model, ordering is preserved.} The
triangle encodes it geometrically. Each child depends only on what the
child above passed down (stage dependency) and each chunk depends only
on the chunk before it at the same child (sequence dependency). Those
two axes -- vertical and horizontal -- are the active constraints. In
this model, the triangle is the tight packing that satisfies both.

This is not only a visualization choice. In this model, the triangle is
the canonical shape of pipelined computation. It is the minimum-area
region in time × stage space that achieves full occupancy while
respecting dependency constraints. Broader regions waste slots; narrower
regions violate ordering.

The triangle is also a \textbf{covering space} (§3.3). The diagonal --
the moment when all children are busy -- is the base space: one ordered
sequence, 1-2-3-4. But each chunk arrived at the diagonal via a
different path through the triangle. Chunk 1 took the longest path
(entered first, fell through all four stages). Chunk 4 took the shortest
(entered last, only at stage 1). Many paths, one output. That is the
covering-map intuition, and in this model it preserves order.

And the triangle is \textbf{fractal}. Zoom into any sub-triangle and you
see the same pattern. If you bundle chunks into mega-chunks, each
mega-chunk moves through a larger triangle the same way a single chunk
moves through a small one. A polysome with 40 ribosomes on an mRNA looks
the same as one ribosome on a short mRNA -- same triangle, different
scale. \emph{Physarum} (slime mold) tendril networks are the same
triangle projected onto geography instead of time × stage.

The top of the triangle has -- one chunk, one path, no parallelism. As
you descend, increases -- more chunks in flight, more independent paths
through the system. At the diagonal, is maximum. Then the ramp-down
triangle on the other side collapses back to zero.

\textbf{Fork is entering the triangle. Race is the diagonal. Fold is
exiting.}

Now zoom out.

The children are standing inside a classroom. The teacher is managing
\emph{three} lines of children, each passing different-colored balls.
When one line stalls -- for example, a dropped ball -- the teacher
slides a waiting child from another line into the gap. No one is idle.
This is \textbf{turbulent multiplexing}: multiple pipelines sharing idle
slots across lines.

Zoom out again. The school is one of many in a district. The district
coordinator doesn't manage individual children or individual balls. She
manages the \emph{shape} of the system -- how many lines, how wide, how
they interconnect. She has discovered that the number of
\emph{independent parallel paths} through the system matters more than
the speed of any individual child. She calls this number .

Zoom out once more. You are looking at a strand of DNA inside a cell,
and the cell is doing a structurally similar thing. The replication fork
is the teacher. Okazaki fragments are the bundled balls. DNA ligase is
the child at the end of the line, stitching fragments together without
requiring global arrival order. The cell has used this pattern for
billions of years.

A working hypothesis follows from this zoom-out: efficient coordination
patterns can be discovered across natural and engineered systems under
shared constraints.

Three natural axioms set the stage.

\begin{itemize}
\tightlist
\item
  \textbf{Locality axiom}: if correctness is governed by local
  constraints, forcing global sequential order adds latency without
  adding truth.\\
\item
  \textbf{Topology axiom}: when multiple paths preserve correctness, a
  high-efficiency policy class is to fork them, race them, then fold
  deterministically.\\
\item
  \textbf{Naturalism axiom}: when the same pattern reappears in
  classrooms, cells and networks, it can indicate a
  substrate-independent computational shape rather than a loose
  metaphor.
\end{itemize}

This paper studies whether the same coordination shape recurs across
these settings. It has \textbf{three} operations: \textbf{fork} work
into parallel paths, \textbf{race} paths against each other,
\textbf{fold} results into a single answer. It has one safety mechanism:
\textbf{vent} -- propagate down, never across. It treats failure as
first-class to minimize wasted work.

These four operations are sufficient to express finite directed acyclic
computation graphs in this paper's mechanized scope. They have a natural
topological characterization in terms of Betti numbers, covering spaces
and homotopy equivalence. Classical queueing theory -- Little's Law,
Erlang's formula and Jackson networks -- appears as the degenerate
one-path case. The quantum-mechanical vocabulary -- superposition,
tunneling, interference, entanglement, collapse -- is used here as
structural correspondence language.

This paper began as a practical problem: a sequential bottleneck in a
distributed inference pipeline. Tokens moved through layer nodes one at
a time, and the obvious optimization -- standard pipelining -- wasn't
good enough. The question was not only ``how do I make the pipeline
faster?'' but ``why is there a pipeline at all?'' That reframing led to
a topology that reduced the measured bottleneck and motivated the
broader framework presented here.

The conveyor belt was a dominant 20th-century abstraction: serialize
work. Fork/race/fold is a correction for workloads with , where
sequential-only structure imposes avoidable coordination cost.

Three bodies of existing theory provided the language for this
correction.

I drew heavily from \textbf{quantum physics}, using selected terms as
computational correspondences: superposition is fork, measurement is
observation, collapse is fold, tunneling is early exit, interference is
consensus, entanglement is shared state (§5). These are structural
mappings, with photosynthetic antenna complexes as a literal quantum
anchor example (§1.5).

The quantum-mechanical vocabulary describes the mapped computational
operations with structural precision in this paper's scope. It is a
modeling lens, not an exclusive language claim.

The second muse is \textbf{fluid dynamics}, whose Reynolds number I
purloin wholesale into computation as the pipeline Reynolds number
(§2.3). Fluid dynamics provides more than vocabulary -- it provides a
useful intuition for \emph{when} fork/race/fold matters in this model.
Just as the Reynolds number predicts when laminar flow becomes
turbulent, indicates when sequential processing should yield to
multiplexed scheduling and when the system begins to lose its laminar
ability to recover from local drops.

The fluid-dynamical framing reveals an inverted scaling property (§2.2):
the worst case is small data, where ramp-up overhead dominates. As data
grows, speedup accelerates toward . In this model, larger workloads can
become favorable once ramp-up overhead is amortized.

The third muse is \textbf{stability theory}. A Foster-Lyapunov drift
schema treats the safe operating region as a small set and asks whether
expected motion points toward it. In this manuscript's geometry, failure
is not ``the bug happened'' but ``the drift field turned outward.'' The
topological deficit is therefore not only an efficiency gap but a
failure coordinate: the gap between the manifold the problem needs and
the one the implementation can actually sustain. An
information-theoretic framing (§6.8) turns the Void into an accounting
ledger: fork creates alternatives, fold compresses to one outcome, and
vented paths carry the bits that cannot be retained

.

\subsection{\texorpdfstring{\textbf{1. Structural Homologies in Natural
Systems}}{1. Structural Homologies in Natural Systems}}\label{structural-homologies-in-natural-systems}

Fork/race/fold is used here as a structural model for natural systems.
The same pattern appears across different substrates. I grade each
mapping:

\begin{itemize}
\tightlist
\item
  \textbf{Grade A}: Quantitative correspondence -- the algorithm's math
  directly models the system with embedded predictive power.\\
\item
  \textbf{Grade B}: Structural homology -- deep structural match,
  genuine design insight, no novel quantitative prediction.
\end{itemize}

Grades are evidentiary tiers, not additive votes. Grade B examples
provide structural context; they are not interchangeable with Grade A
quantitative confirmations.

In this framing, the novelty is making the convergence explicit and
testable.

In this paper, computational beauty is treated as an interpretive lens
tied to low topological deficit, explicit invariants, and reproducible
gains under stated assumptions.

\subsubsection{\texorpdfstring{\textbf{1.1 \emph{Physarum polycephalum}:
The Phineas Gage of Distributed Intelligence (Grade
A)}}{1.1 Physarum polycephalum: The Phineas Gage of Distributed Intelligence (Grade A)}}\label{physarum-polycephalum-the-phineas-gage-of-distributed-intelligence-grade-a}

In 1848, a railroad construction foreman named Phineas Gage survived a
43-inch iron rod through his left frontal lobe. He could walk, talk, and
reason -- but his personality changed utterly. The accident revealed
that personality \emph{lives somewhere specific} in the brain. It was
the founding observation of neuropsychology: a single dramatic injury
that reorganized an entire field's understanding of where intelligence
resides.

His skull and scientific divining rod are now in a display case in a
medical library at Harvard.

As Phineas Gage was to neuroscience, so slime mold is to distributed
intelligence.

In 2010, Tero et al.~placed oat flakes on a wet surface in positions
corresponding to the 36 stations of the greater Tokyo rail network

. They introduced a single \emph{Physarum polycephalum} slime mold at
the position corresponding to Tokyo station. The organism -- which has
no brain, no neurons, no central nervous system of any kind -- extended
exploratory tendrils in all directions (\textbf{fork}). Multiple
tendrils reached each food source via different routes (\textbf{race}).
The organism then pruned inefficient connections, reinforcing high-flow
tubes and abandoning low-flow ones (\textbf{fold} with \textbf{venting}
of abandoned paths).

Within 26 hours, the slime mold had independently constructed a
transport network with tradeoffs similar to the actual Tokyo rail system
-- a network that professional engineers had spent decades and billions
of dollars optimizing

.

The correspondence Tero et al.~emphasize is at the level of system
tradeoffs:

\begin{itemize}
\tightlist
\item
  \textbf{Cost}: the network remained competitive on total length\\
\item
  \textbf{Fault tolerance}: it retained resilience under link removal\\
\item
  \textbf{Transport efficiency}: it balanced transport performance
  against cost in the same design space\\
\item
  \textbf{Topology}: it remained cyclic rather than collapsing to a
  single path
\end{itemize}

The author does not know of a slime mold yet on display at a public
library, but remains hopeful.

The mapping to fork/race/fold is presented as an operational mechanism
in this model:

\textbf{Physarum Behavior} \textbar{} \textbf{Fork/Race/Fold Operation}
\textbar{}

Exploratory tendril extension \textbar{} \textbf{Fork}: create parallel
paths from current position \textbar{}

Cytoplasmic streaming through tubes \textbar{} \textbf{Race}: flow rate
determines winner \textbar{}

Tube reinforcement (positive feedback) \textbar{} \textbf{Fold}:
high-flow paths become canonical \textbar{}

Tube abandonment (starvation) \textbar{} \textbf{Vent}: low-flow paths
released, descendants shed \textbar{}

Shuttle streaming (oscillatory flow) \textbar{} \textbf{Self-describing
frames}: bidirectional flow carries positional information \textbar{}

Just as Phineas Gage's injury revealed that intelligence has a specific
anatomical locus, \emph{Physarum}'s rail network shows that optimization
can emerge without centralized cognition. Fork/race/fold does not
require a neural substrate; it needs parallel paths, a selection signal
and a way to prune. In this manuscript, that structure is instantiated
in protoplasm, silicon and 10-byte frame transport.

\textbf{Predictive power}: The Wallington Rotation's chunk-size framing
predicts a cubic tradeoff between carried flow and tube radius. That is
qualitatively consistent with Murray's-law analyses of \emph{Physarum}
tube morphology

.

\subsubsection{\texorpdfstring{\textbf{1.2 DNA Replication: The Original
Self-Describing Frame Protocol (Grade
A)}}{1.2 DNA Replication: The Original Self-Describing Frame Protocol (Grade A)}}\label{dna-replication-the-original-self-describing-frame-protocol-grade-a}

DNA's two strands run antiparallel. The leading strand synthesizes
continuously (clean pipeline). The lagging strand produces
\textbf{Okazaki fragments} -- 1,000--2,000 nucleotide chunks in
prokaryotes, 100--200 in eukaryotes -- synthesized out of order and
stitched together by DNA ligase

.

Each Okazaki fragment is a \textbf{self-describing frame}: its genomic
coordinate is its stream\_id + sequence. DNA ligase is the \textbf{frame
reassembler} -- it joins fragments without requiring global ordering.
The replication fork moves at \textasciitilde1,000 nt/s in \emph{E.
coli}. At any moment, 1--3 fragments are being synthesized
simultaneously, giving a pipeline Reynolds number --.

\textbf{Predictive power}: My chunked pipeline formula predicts that
prokaryotic fragments (\textasciitilde1,000 nt) should be longer than
eukaryotic fragments (\textasciitilde150 nt) because eukaryotes have
more processing stages (chromatin reassembly, histone deposition). This
is directionally consistent with reported fragment ranges. The framework
also predicts that organisms with lower (more exposed single-stranded
DNA during lagging strand synthesis) should show stronger strand
asymmetries. Directionally similar asymmetries are well documented in
bacterial genomes

.

\subsubsection{\texorpdfstring{\textbf{1.3 Saltatory Nerve Conduction:
The Formula Tracks Measured Range (Grade
A)}}{1.3 Saltatory Nerve Conduction: The Formula Tracks Measured Range (Grade A)}}\label{saltatory-nerve-conduction-the-formula-tracks-measured-range-grade-a}

In myelinated neurons, action potentials jump between nodes of Ranvier
(\textasciitilde1--2 mm apart) instead of propagating continuously.
Multiple action potentials are in-flight simultaneously across different
internodal segments.

Perhaps you picture biological denial of service-style ``packet
overload''? The biology includes refractory dynamics that suppress
re-firing for a short window after a node activates. This acts like a
one-way valve/buffer so that, at the modeled level, in-flight action
potentials remain directionally separated.

This is chunked pipelining. The ``chunking'' allows the brain to receive
a high-frequency stream of data rather than a single pulse. It allows
for more nuanced signaling---the frequency of the spikes (the
``bitrate'') conveys the intensity of the stimulus.

By only depolarizing the membrane at the nodes of Ranvier, the neuron
saves a massive amount of metabolic energy (ATP) that would otherwise be
spent pumping ions back and forth across the entire length of the axon.
The brain doesn't just pipeline; it optimizes for energy efficiency by
reducing the number of times ions need to be pumped back across the
membrane.

The Wallington formula gives the right order of magnitude for saltatory
conduction when stage length is interpreted as internodal distance and
stage delay as nodal regeneration time:

Using representative millimetric internodes and rapid nodal regeneration
gives a velocity on the order of m/s, which matches the measured range
for large myelinated fibers. Classical work established the
electro-saltatory nature of nodal regeneration

, later experiments showed that increasing internodal distance
accelerates conduction toward a broad plateau

, and review literature summarizes measured conduction velocities of
roughly 10-120 m/s in large myelinated fibers and about 0.5-2 m/s in
unmyelinated fibers

. The point here is order-of-magnitude agreement, not an exact
one-parameter biological law.

\textbf{Measured conduction velocity: about 100 m/s} (for large
myelinated fibers). Unmyelinated conduction is typically much slower,
around 0.5-2 m/s

. The large speedup is real, measured, and directionally consistent with
the same stage-length / stage-delay intuition that predicts my pipeline
speedups.

Myelin is the biological argument for investing in transport-layer
reliability to enable larger chunks -- skip intermediate processing,
insulate the wire. This is the case for UDP over TCP: invest in framing
reliability so you can skip ordered delivery.

\subsubsection{\texorpdfstring{\textbf{1.4 Polysome Translation: The
Wallington Rotation in Biology (Grade
B)}}{1.4 Polysome Translation: The Wallington Rotation in Biology (Grade B)}}\label{polysome-translation-the-wallington-rotation-in-biology-grade-b}

A so-called polysome consists of multiple ribosomes simultaneously
translating the same mRNA, spaced \textasciitilde30--40 codons apart.
This \emph{is} the Wallington Rotation: the mRNA is the pipeline, each
ribosome processes a chunk, and multiple proteins emerge concurrently.

A back-of-envelope throughput illustration using nominal translation
rates gives: without pipelining, 40 proteins from one mRNA is
\textasciitilde2,400 s; with polysome overlap, \textasciitilde118 s
(about 20x). This is an illustrative order-of-magnitude comparison, not
a calibrated cross-lab benchmark.

When drops below \textasciitilde0.6, the mRNA is targeted for
degradation (no-go decay). The cell destroys underutilized pipelines and
reallocates ribosomes -- behavior that qualitatively aligns with
turbulent-multiplexing intuition. Under stress, cells globally reduce
but maintain high on priority mRNAs via IRES elements.

\subsubsection{\texorpdfstring{\textbf{1.5 Photosynthetic
Light-Harvesting: Fork/Race at Quantum Scale (Grade
A)}}{1.5 Photosynthetic Light-Harvesting: Fork/Race at Quantum Scale (Grade A)}}\label{photosynthetic-light-harvesting-forkrace-at-quantum-scale-grade-a}

Photosynthetic light-harvesting is widely considered a strong example of
quantum-coherent energy transfer in biology. In these systems, the
algorithm in action is environment-assisted quantum transport, where
excitons exploit spatial superposition to sample multiple pathways and
achieve high reported transfer efficiency before decoherence.

Antenna complexes in photosynthesis contain large pigment networks.
Photon excitation energy forks across the pigment network, races through
multiple pathways and the first path to reach the reaction center wins.
Charge separation is fold. Non-photochemical quenching is venting. At
this step, excitation-energy transfer is exceptionally efficient and is
often described as near-unity relative to the much lower whole-system
photosynthetic yield once downstream losses are included

. The fork/race/fold characterization applies to the excitation-transfer
step specifically.

Fleming et al.~(2007) showed long-lived quantum-coherent signatures in
photosynthetic energy transfer

. The fork/race/fold framing predicts that transfer efficiency should
increase with pigment count (more forked paths = higher probability of
reaching the reaction center before decoherence) but with diminishing
returns once reaction-center capture is already highly efficient. The
quantum vocabulary in §5 is used as structural correspondence language.

\subsubsection{\texorpdfstring{\textbf{1.6 Immune System V(D)J
Recombination (Grade
B)}}{1.6 Immune System V(D)J Recombination (Grade B)}}\label{immune-system-vdj-recombination-grade-b}

The adaptive immune system generates unique antibody configurations
through combinatorial recombination (\textbf{fork}), exposes them to
antigen simultaneously (\textbf{race}) and expands the winners through
clonal selection (\textbf{fold}). Non-binding clones are eliminated
(\textbf{vent}). Self-reactive B cells undergo clonal deletion -- the
lineage is eliminated, but sibling B cells with different recombinations
are unaffected. The implied parallelism factor is on the order of .

This is not just parallelism; it is \textbf{probabilistic parallelism}.
The immune system does not know which configuration will bind the
antigen. It forks a vast library, races them against the antigen and
folds the winners. The structure is strongly analogous to the
fork/race/fold pattern used in distributed systems.

\subsubsection{\texorpdfstring{\textbf{1.7 Transformers: Accidental
Rediscovery (Grade
A)}}{1.7 Transformers: Accidental Rediscovery (Grade A)}}\label{transformers-accidental-rediscovery-grade-a}

The biological examples above are evolutionary discoveries. But the
pattern extends to human-engineered systems that arrived at similar
structure without topological framing. The transformer architecture
(§6.11) is the clearest example: multi-head attention is fork/race/fold
( heads fork, compute in parallel, concatenate-and-project to fold),
feed-forward layers are fork/vent/fold (expand 4x, activate/suppress,
contract), residual connections are two-path fork with additive fold,
and softmax is continuous venting. The architecture can be represented
as a nested fork/race/fold graph.

\textbf{Falsifiable prediction.} If the fork/race/fold characterization
is structural rather than coincidental, then removing parallel-path
properties while preserving parameter count should degrade performance
more than a naive parameter-count argument would suggest. Voita et
al.~(2019)

showed that attention heads are unequally specialized and that pruning
has non-uniform quality impact, with the most important heads often
associated with syntactic functions. That is consistent with, but does
not prove, the broader topological hypothesis advanced here. Similarly,
replacing the 4x FFN expansion-contraction (fork/vent/fold) with a
same-parameter linear layer should degrade performance if activation
sparsity is doing real path selection work. The prediction is testable:
ablate the parallel-path structure, measure of the modified
architecture, and correlate that with the performance drop. A
deterministic companion toy-attention harness now makes the narrow
version executable at fixed parameters
(companion-tests/artifacts/toy-attention-fold-ablation.\{json,md\}):
keeping keys, values, score function and query grid fixed, linear fold
reconstructs the teacher exactly (MSE 0.000), winner-take-all incurs MSE
0.163, and early-stop incurs MSE 1.071; bootstrap 95\% intervals over
the query grid are {[}0.000, 0.000{]}, {[}0.120, 0.211{]}, and {[}0.858,
1.288{]} respectively, and the exact-within-0.01 rates are 1.000, 0.185,
and 0.074. A seeded Gnosis cancellation benchmark then lifts the same
claim into a learned setting
(companion-tests/artifacts/gnosis-fold-training-benchmark.\{json,md\}):
three parameter-matched .gg programs keep topology, parameter count, and
data fixed and differ only in FOLD \{ strategy: \ldots{} \}. Across 8
seeds on the cancellation-sensitive target family left - right, linear
fold reaches eval MSE 0.000 with 95\% seed-bootstrap interval {[}0.000,
0.000{]} and exact-within-0.05 rate 1.000, winner-take-all reaches 0.408
with interval {[}0.396, 0.421{]} and exact rate 0.038, and early-stop
reaches 0.735 with interval {[}0.732, 0.740{]} and exact rate 0.000;
cancellation-line mean absolute error is 0.000, 0.834, and 0.764
respectively. A harder seeded Gnosis mini-MoE routing benchmark
(companion-tests/artifacts/gnosis-moe-routing-benchmark.\{json,md\})
keeps a four-expert routed topology and a fixed 16-parameter budget
while swapping only the fold rule: linear fold reaches eval MSE 0.001
with 95\% seed-bootstrap interval {[}0.001, 0.001{]}, winner-take-all
reaches 0.328 with interval {[}0.267, 0.389{]}, and early-stop reaches
0.449 with interval {[}0.444, 0.457{]}; the dual-active-region absolute
error is 0.027, 0.402, and 0.474, showing the same ranking on a
genuinely routed learned task where two paths should contribute. A
paired negative-control artifact
(companion-tests/artifacts/gnosis-negative-controls.\{json,md\}) then
reuses those same affine and routed topologies on one-path task families
where additive recombination is unnecessary; there the separation
collapses exactly as predicted, with all three fold rules reaching eval
MSE 0.000, exact-within-tolerance rate 1.000, and zero inter-strategy
gap on both the affine-left-only and positive-x single-expert controls.
The next layer is now a regime map rather than just endpoint witnesses
(companion-tests/artifacts/gnosis-fold-boundary-regime-sweep.\{json,md\}):
in the affine family the first separated regime appears at cancellation
weight 0.500 and the final linear-vs-best-nonlinear eval-MSE advantage
reaches 0.393, while in the routed family the first separated regime
appears at dual-activation weight 0.750 and the final linear advantage
reaches 0.111. The converse side is also executable in
companion-tests/artifacts/gnosis-adversarial-controls-benchmark.\{json,md\}:
on the winner-aligned signed max-by-magnitude affine task,
winner-take-all beats linear and early-stop in final eval MSE (0.095 vs
0.271 vs 0.484) and learning-curve area (0.097 vs 0.265 vs 0.482), while
on the short-budget x-priority routed task and left-priority affine
task, early-stop has the best learning-curve area (0.036 and 0.006)
against linear (0.200 and 0.531) and winner-take-all (0.342 and 1.375).
So the evidence boundary is symmetric: linear wins when cancellation or
dual contribution is truly required; sparse selection wins when the
target family itself is sparse and ordered.

\subsubsection{\texorpdfstring{\textbf{1.8 Observed Recurrence in
Selected
Examples}}{1.8 Observed Recurrence in Selected Examples}}\label{observed-recurrence-in-selected-examples}

In this manuscript's selected examples -- spanning roughly seven orders
of magnitude from quantum-coherent pigment networks to billion-parameter
neural networks -- different substrates exhibit a related structural
motif. This is suggestive rather than conclusive, and is used here as a
bounded correspondence claim under three constraints:

\begin{enumerate}
\def\labelenumi{\arabic{enumi}.}
\tightlist
\item
  \textbf{Finite resources, high demand} → chunked pipelining and
  multiplexing\\
\item
  \textbf{Unknown correct answer} → fork/race/fold with vent\\
\item
  \textbf{No global clock} → self-describing frames with out-of-order
  reassembly
\end{enumerate}

These three constraints make fork/race/fold a strong candidate for high
efficiency within the class of finite DAG topologies examined in this
paper. Systems lacking all three can still use fork/race/fold
(transformers have synchronized SGD; photosynthesis has electromagnetic
field synchronization). In the finite constructions evaluated here, when
all three constraints bind simultaneously, no outperforming alternative
was observed in the tested topology set on the measured criteria. Other
concurrent models (gossip protocols, epidemic algorithms,
eventually-consistent CRDTs) also operate under these constraints but
make different tradeoffs --- gossip sacrifices deterministic fold for
probabilistic convergence, CRDTs sacrifice fold entirely for
commutativity. The claim is not unique optimality; it is selective
evidence for pressure toward this topology when systems require both
parallelism and deterministic reconciliation. The conveyor belt remains
the degenerate one-path case.

\subsection{\texorpdfstring{\textbf{2. The
Algorithm}}{2. The Algorithm}}\label{the-algorithm}

\subsubsection{\texorpdfstring{\textbf{2.1 Pipeline
Model}}{2.1 Pipeline Model}}\label{pipeline-model}

A pipeline with stages processes workload of items.

\textbf{Serialized}: -- each item completes all stages before the next
begins.

\textbf{Pipelined}: -- stage-local ordering suffices.

\textbf{Chunked}: -- intra-stage parallelism (SIMD, batched ops) with
chunk size .

For chunks across stages, the idle fraction during ramp-up/ramp-down:

\subsubsection{\texorpdfstring{\textbf{2.2 The Inverted Scaling
Property}}{2.2 The Inverted Scaling Property}}\label{the-inverted-scaling-property}

Under the idealized scheduling model used for this derivation, two
assumptions are explicit: (A1) per-chunk stage service times are
homogeneous across stages, and (A2) inter-stage
communication/synchronization cost is zero. Under A1-A2, the speedup of
chunked pipelining over serialized processing is:

This framing is complementary to classical parallel-scaling laws:
Amdahl's fixed-workload limit and Gustafson's scaled-workload
reformulation

. Here, is used as a structural diagnostic for where serial fractions
are imposed by topology, not as a replacement for those bounds.

For small (few items, few stages), the denominator's term -- the ramp-up
cost -- dominates. The pipeline spends most of its time filling and
draining, never reaching full occupancy. The idle fraction is large.
This is the \textbf{worst case}.

For large , the term dominates and becomes negligible:

Under A1-A2, the speedup approaches -- the product of chunk size and
stage count. The pipeline is fully occupied. The kids all have balls.
Idle fraction approaches zero. The kids are all juggling.

The technique gets \emph{faster and faster} as the work at hand grows.

This is a profoundly inverted scaling property, which is a useful and
unusual feature of the algorithm. In most engineering contexts, the hard
problem is scale -- systems that work beautifully on small inputs
collapse under large ones. Here, the opposite is true: large datasets
are where pipelining shines, approaching its theoretical maximum speedup
of . Small datasets are where the overhead hurts.

The optimization challenge in fork/race/fold-based pipelines is not
``how do I survive at scale?'' but ``how do I avoid overpaying on small
workloads?'' -- a far more pleasant problem. Trivial solutions for such
trivial concerns are panoply -- early exit, dynamic chunk sizing and
adaptive scheduling, for example -- but speedups are more precious when
they arrive in abundance.

The fluid-dynamical analogy (§2.3) captures this behavior. Low (many
chunks, few stages -- large data) corresponds to laminar flow: smooth,
predictable, high utilization. High (few chunks, many stages -- small
data) corresponds to turbulent flow: idle slots appear, multiplexing
becomes necessary, overhead rises. The Reynolds number indicates the
crossover region, and the laminar regime is the one that grows with data
size.

\subsubsection{\texorpdfstring{\textbf{2.3 The Pipeline Reynolds Number
and Failure
Turbulence}}{2.3 The Pipeline Reynolds Number and Failure Turbulence}}\label{the-pipeline-reynolds-number-and-failure-turbulence}

I define:

This is the ratio of stages to chunks -- the density of the pipeline.
Low (): laminar regime, steady-state, high utilization, and strong
recovery after local perturbation. Transitional (--): idle-slot recovery
is profitable and the system benefits from explicit restorative
pressure. High (): turbulent regime, multiplexing across requests yields
the largest benefit, and failure begins to behave like a geometric
instability rather than a rare event.

The Reynolds-number mapping is an explicit analogy. In fluid dynamics,
predicts transition from laminar to turbulent flow: inertial forces
(numerator) versus viscous forces (denominator). In computation, the
correspondence used here is: stages as inertial pressure (more stages =
more work in flight), and chunks as viscous pressure (larger chunks =
more resistance to context switching). Low (large chunks, few stages) is
laminar-like; high (small chunks, many stages) is turbulent-like. The
transition occurs when ramp-up/ramp-down idle cost exceeds multiplexing
recovery benefit.

The geometric interpretation is a drift test. Let be arrival pressure,
be service pressure, and be restorative vent pressure (reneging,
shedding, or other fail-flow relief). Then the local stability margin
is:

When , the transition kernel still points inward toward the small set of
safe states. When , the pipeline has lost its restorative gravity:
backlog expands, the manifold flattens, and failure becomes outward
drift rather than a binary fault. High is therefore not merely idle
slots; it is the onset of \textbf{failure turbulence}.

\subsubsection{\texorpdfstring{\textbf{2.4 Four
Primitives}}{2.4 Four Primitives}}\label{four-primitives}

Given pipeline state and operation set :

\begin{enumerate}
\def\labelenumi{\arabic{enumi}.}
\tightlist
\item
  \textbf{Fork}: -- create independent branch states, each processing a
  subset of . Topological effect: .\\
\item
  \textbf{Race}: -- advance all branches concurrently; select the first
  to reach a valid completion. Losers are vented. Exploits homotopy
  equivalence.\\
\item
  \textbf{Fold}: -- wait for all branches to complete (or vent); apply
  deterministic merger to produce a single canonical state. Topological
  effect: .\\
\item
  \textbf{Vent}: -- cease output, recursively vent all descendants,
  leave siblings untouched. \textbf{One rule: propagate down, never
  across.} The system releases excess energy from paths that cannot
  contribute useful work -- a pressure relief valve that prevents
  computational overheating.
\end{enumerate}

\textbf{Completeness (finite, mechanized scope).} These four primitives
are sufficient to express finite directed acyclic computation graphs
under explicit decomposition assumptions. Any finite DAG can be
decomposed into fork points (nodes with out-degree \textgreater{} 1),
join points (nodes with in-degree \textgreater{} 1) and linear chains.
Fork creates divergences. Fold creates convergences. Race is fold with
early termination. Vent handles failures and excess energy. Linear
chains are the trivial case (no fork, no fold). In the formal stack,
local decomposition is constructive and the global statement is an
explicit-assumption theorem schema, paired with executable finite-DAG
decomposition checks

.

\subsubsection{\texorpdfstring{\textbf{2.5 Correctness
Conditions}}{2.5 Correctness Conditions}}\label{correctness-conditions}

Fork/race/fold preserves correctness when:

\begin{itemize}
\tightlist
\item
  \textbf{C1 (Constraint locality)}: Stage-local ordering is sufficient
  for global correctness.\\
\item
  \textbf{C2 (Branch isolation)}: A vented branch does not corrupt
  siblings.\\
\item
  \textbf{C3 (Deterministic fold)}: The merger is deterministic.\\
\item
  \textbf{C4 (Termination)}: Every branch either completes, is vented,
  or times out in finite time.
\end{itemize}

These conditions are mechanized in a two-layer formal stack.
Finite-state models in TLA+ verify C1--C4 as invariants across the
formal module set (checked by both TLC and the self-hosted aeon-logic
parser/checker). Lean 4 theorem schemas verify the quantitative
identities that depend on C1--C4 under explicit assumptions. The
sufficiency claim --- that any finite DAG decomposes into fork points,
join points, and linear chains, and that these four conditions preserve
correctness through the decomposition --- is verified constructively by
executable finite-DAG decomposition checks

.

\subsubsection{\texorpdfstring{\textbf{2.6 Five Fold
Strategies}}{2.6 Five Fold Strategies}}\label{five-fold-strategies}

Not all folds are equal. The choice of merger determines the
computational semantics:

\textbf{Strategy} \textbar{} \textbf{Semantics} \textbar{} \textbf{Time
Complexity} \textbar{} \textbf{When} \textbar{}

\textbf{Winner-take-all} \textbar{} Best result by selector \textbar{}
comparisons \textbar{} One answer needed, clear criterion \textbar{}

\textbf{Quorum} \textbar{} of must agree \textbar{} pairwise comparisons
\textbar{} Byzantine fault tolerance \textbar{}

\textbf{Merge-all} \textbar{} All results contribute \textbar{} where =
merger cost \textbar{} Complementary information \textbar{}

\textbf{Consensus} \textbar{} Constructive/destructive interference
\textbar{} pairwise comparisons \textbar{} Signal amplification or
outlier detection \textbar{}

\textbf{Weighted} \textbar{} Authority-weighted merger \textbar{} where
= merger cost \textbar{} Heterogeneous source quality \textbar{}

Race is not a fold strategy -- it is a separate primitive. Race picks
the \emph{fastest} result. Winner-take-all picks the \emph{best} result.
The distinction matters: race terminates early (venting losers),
winner-take-all waits for all branches to complete.

Implicit in this is the fact that failure is a necessary component of
any robust system. Failure modes are handled by the vent primitive,
which propagates down the tree but never across branches. This ensures
that a failure in one branch does not cascade to other branches,
maintaining the isolation property required for correctness. A system
that cannot fail gracefully is not robust.

\subsubsection{\texorpdfstring{\textbf{2.7 Vent
Propagation}}{2.7 Vent Propagation}}\label{vent-propagation}

Venting is the protocol-level analogue of NaN propagation in IEEE 754,
AbortSignal in web APIs and apoptosis in biology. The one rule --
\textbf{propagate down, never across} -- makes composition safety an
architectural feature rather than an accidental one. Under C2 (branch
isolation), fork/race/fold compositions preserve this safety property
because venting never crosses branch boundaries.

\subsubsection{\texorpdfstring{\textbf{2.8 The Worthington
Whip}}{2.8 The Worthington Whip}}\label{the-worthington-whip}

The Worthington Whip extends fold for aggressive parallel shard merging.
A single workload of items is sharded across parallel pipelines, each
processing items. At fold, a cross-shard correction reconciles the
results. In the geometric framing, this reconciliation is
\textbf{tension} rather than incidental cleanup: the manifold
accumulated debt while it widened, and the whip snap is the controlled
payment that returns it to a single stable sheet.

\textbf{Derivation of the} \textbf{reduction.} In a computation with
pairwise dependencies, an unsharded system processes all pairs. After
sharding into partitions of items each, each shard processes only its
intra-shard pairs: . The total intra-shard work across all shards is .
The cross-shard pairs --- the ones not processed within any shard ---
number . As , the ratio of intra-shard work to total work approaches ,
so the per-shard compute reduction is . Per shard, each shard avoids of
the total pairs, but since each shard processes of the total, the
per-shard savings relative to processing the full is . The cross-shard
correction at fold time reconciles the missing pairs --- this is the
whip snap.

The fold phase is the whip snap: all parallel shards converge to a
single definite state. The computational snap is a single-state
reconciliation step that can preserve substantial parallel gains while
re-entering a canonical sequential state. The important point is that
collapse is never free. The snap is the visible moment at which
reconciliation debt is paid.

\textbf{Beyond pairwise:} \textbf{-tuple dependencies.} The derivation
assumes pairwise interactions. For -tuple dependencies (e.g., 3-way
consistency checks), the intra-shard work fraction generalizes to ,
which approaches as . The per-shard savings grow with : for , the
reduction is --- sharding is typically more beneficial for higher-order
dependencies. The cross-shard correction at fold reconciles missing
-tuples. For unequal partition sizes, the correction cost depends on the
size distribution; the formula is the symmetric optimum, and partition
skew increases the cross-shard fraction.

\textbf{The crossover point.} Adding shards reduces per-shard
computation by but increases cross-shard correction cost. The correction
is a fold over partial results --- itself an operation for merge-all or
for quorum. The crossover occurs when the marginal correction cost of
shard exceeds the marginal per-shard savings: . For pairwise
dependencies with merge-all fold, this yields the model optimum . The
TLA+ WhipCrossover model mechanizes this in bounded finite-state space:
it explores configurations up to explicit bounds and verifies the
inequality within that scope

.

These whipper-snapper folds are an aggressive expression of parallel
shard reconciliation in this framework.

\subsection{\texorpdfstring{\textbf{3. The Topology of
Fork/Race/Fold}}{3. The Topology of Fork/Race/Fold}}\label{the-topology-of-forkracefold}

\subsubsection{\texorpdfstring{\textbf{3.1 Betti Numbers Classify
Computation
Graphs}}{3.1 Betti Numbers Classify Computation Graphs}}\label{betti-numbers-classify-computation-graphs}

The first Betti number counts independent parallel paths in a
topological space. In this framework, it is a primary control variable:

\textbf{Structure} \textbar{} \textbf{β1\hspace{0pt}} \textbar{}
\textbf{Parallelism} \textbar{} \textbf{Fault Tolerance} \textbar{}

Sequential pipeline \textbar{} 0 \textbar{} None \textbar{} None
\textbar{}

Single fork/join \textbar{} 1 \textbar{} One level \textbar{} One
failure \textbar{}

Fork with paths \textbar{} \textbar{} -way \textbar{} \$K-1 failures
\textbar{}

Full mesh of nodes \textbar{} \textbar{} Maximum \textbar{} Maximum
\textbar{}

Fork/race/fold is the operation that \textbf{temporarily raises}
\textbf{to exploit parallelism, then lowers it back to zero}:

\textbf{Primitive} \textbar{} \textbf{Topological Operation} \textbar{}
\textbf{Effect on β1\hspace{0pt}} \textbar{}

\textbf{Fork} \textbar{} Create parallel paths \textbar{} \textbar{}

\textbf{Race} \textbar{} Traverse homotopy-equivalent paths \textbar{}
stays high \textbar{}

\textbf{Fold} \textbar{} Merge all paths to single output \textbar{}
\textbar{}

\textbf{Vent} \textbar{} Release a path \textbar{} \textbar{}

Many historical process designs -- Ford's assembly line, TCP's ordered
byte stream, hospital referral chains, T+2 financial settlement -- can
be interpreted as forcing onto problems whose natural topology has .
Healthcare diagnosis has intrinsic (blood work, imaging, genetic
screening, and specialist consultation are independent). The referral
system forces . The mismatch correlates with multi-year diagnostic
delay: the 2024 EURORDIS Rare Barometer diagnosis survey reports an
average diagnosis time of 5 years for people living with a rare disease

. Financial settlement has intrinsic . T+2 forces . Using the DTCC/NSCC
2024 average daily transaction value baseline of \$2.219 trillion

, a simple two-day lockup heuristic implies on the order of \$4.4
trillion tied up during T+2 settlement; larger figures discussed in the
companion suite are model outputs rather than DTCC-reported statistics

.

\subsubsection{\texorpdfstring{\textbf{3.2 Homotopy
Equivalence}}{3.2 Homotopy Equivalence}}\label{homotopy-equivalence}

Two computations are homotopy equivalent if they produce the same result
through different topological paths. In a sequential pipeline, there is
exactly one path -- no homotopy is possible. In a fork/race graph with
paths, if the computation is deterministic, all paths are homotopy
equivalent.

\textbf{Race exploits homotopy equivalence}: race discovers that all
paths lead to the same answer and takes the fastest. \textbf{Fold
handles the general case}: when paths are \emph{not} homotopy equivalent
(a blood test and an MRI give different information), the merger
function combines non-equivalent results into a richer output than any
single path could provide.

The distinction is topological: race requires homotopy equivalence
(-trivial computation on each path). Fold does not. This is why they are
separate primitives.

\subsubsection{\texorpdfstring{\textbf{3.3 Covering Spaces and
Self-Describing
Frames}}{3.3 Covering Spaces and Self-Describing Frames}}\label{covering-spaces-and-self-describing-frames}

A covering space maps onto a base space such that every point has a
neighborhood that is evenly covered. Self-describing frames create a
covering space over the computation graph. Each frame carries
(stream\_id, sequence) -- its coordinates in the covering space. The
base space is the sequential computation. The covering space is the
multiplexed computation.

The \textbf{frame reassembler is the covering map}: it projects the
cover back to the base space. Frames arrive from any point in the cover
(any stream, any sequence) and are reassembled into sequential order.

\textbf{TCP primarily exposes a base-space abstraction} -- one ordered
byte stream. Simply connected. \textbf{UDP with self-describing frames
exposes a covering-space abstraction} -- many streams, local ordering,
out-of-order reassembly. The topological degree of the covering map is
the multiplexing factor.

This is precisely what DNA ligase does: Okazaki fragments arrive from
the covering space (out-of-order lagging-strand synthesis) and are
projected back to the base space (the complete double-stranded genome).
DNA ligase is the covering map. It has been performing this topological
operation for 4 billion years.

\subsubsection{\texorpdfstring{\textbf{3.4 The Fundamental Group and
Protocol
Design}}{3.4 The Fundamental Group and Protocol Design}}\label{the-fundamental-group-and-protocol-design}

The fundamental group classifies loops up to homotopy:

\begin{itemize}
\tightlist
\item
  \textbf{TCP}: . One path. Simply connected. Works for simply connected
  problems.\\
\item
  \textbf{HTTP/2}: Application layer has (multiplexed streams), but TCP
  substrate has (one ordered byte stream). \textbf{This is a topological
  mismatch.} Head-of-line blocking is the symptom: losing one packet on
  any stream blocks \emph{all} streams because the underlying space
  cannot support independent paths.\\
\item
  \textbf{HTTP/3 (QUIC)}: Partially resolves the contradiction with
  per-stream independence on UDP. But maintains ordered delivery within
  each stream -- within each stream is trivial.\\
\item
  \textbf{Aeon Flow over UDP}: Self-describing frames in the covering
  space. No ordered delivery anywhere. of the wire is designed to match
  of the application. This removes ordered-delivery coupling as a
  head-of-line source in the modeled transport stack.
\end{itemize}

\subsubsection{\texorpdfstring{\textbf{3.5 Time-Indexed Topological
Filtration}}{3.5 Time-Indexed Topological Filtration}}\label{time-indexed-topological-filtration}

The evolution of over a computation's lifetime forms a \emph{filtration}
--- a nested sequence of topological spaces indexed by time:

\begin{itemize}
\tightlist
\item
  : Computation starts. .\\
\item
  : jumps to .\\
\item
  During race: stays at .\\
\item
  : drops by 1 per vented path.\\
\item
  : .
\end{itemize}

\textbf{Terminology note.} This time-indexed evolution is a
\emph{filtration} in the algebraic-topological sense --- a nested
sequence of subcomplexes where each is the computation graph at time .
It borrows the birth/death language of persistent homology (Edelsbrunner
et al., 2002

) but the filtration parameter is \emph{time}, not \emph{distance scale}
as in classical topological data analysis of point clouds. The formal
properties of TDA persistence (stability under perturbation, isometry
invariance) apply only when the filtration is metric-indexed; the
time-indexed version retains the birth/death structure but not the
stability guarantees.

The filtration diagram encodes: how much parallelism was used (features
born at fork), how quickly bad paths were pruned (short persistence =
speculation), how much redundancy survived to fold (long persistence =
consensus). A well-optimized system has short vent persistence (release
early) and long fold persistence (exploit parallelism fully).

\subsubsection{\texorpdfstring{\textbf{3.6 Category-Theoretic
Framing}}{3.6 Category-Theoretic Framing}}\label{category-theoretic-framing}

In category theory, a so-called monoidal category is a mathematical
system consisting of a collection of objects and morphisms, or a way to
combine objects in a way similar to multiplication.

Fork/race/fold forms a \textbf{monoidal category}:

\begin{itemize}
\tightlist
\item
  \textbf{Objects}: computation states (sets of active streams).\\
\item
  \textbf{Morphisms}: Fork (), Race (), Fold ().\\
\item
  \textbf{Tensor product} : parallel composition.\\
\item
  \textbf{Composition} : sequential composition.
\end{itemize}

The conveyor belt uses only composition. Fork/race/fold uses both
composition and tensor product. In this sketch, that suggests a broader
expressive surface. Vent propagation is modeled as a \textbf{natural
transformation} from active computations to terminated computations --
preserving morphism structure across the tensor product, i.e.,
``propagate down, never across.''

\textbf{What this buys.} The monoidal framing is not fully developed
here --- a rigorous treatment would require specifying the unit object
(the empty computation), proving associativity and unit laws for both
and , and verifying the interchange law. The claim is sketched, not
proved, and the paper's results do not depend on it. The value is
taxonomic: it connects fork/race/fold to the category-theoretic
literature on dataflow (hylomorphisms = unfold/fold, which map to
fork/fold) and allows future work to import results from monoidal
category theory --- in particular, the coherence theorem (Mac Lane, 1963

) would guarantee that different compositions of fork/race/fold reach
the same result regardless of bracketing, which is a stronger form of C3
(deterministic fold).

\subsection{\texorpdfstring{\textbf{4. Containing Queueing
Theory}}{4. Containing Queueing Theory}}\label{containing-queueing-theory}

\subsubsection{\texorpdfstring{\textbf{4.1 Little's Law as a Special
Case}}{4.1 Little's Law as a Special Case}}\label{littles-law-as-a-special-case}

Little's Law states: , where is the average number of items in a system,
is the arrival rate and is the average time in the system. This is the
foundational result of queueing theory, proved by Little in 1961

and considered universal within its domain.

\textbf{Containment result (operational form).} Under assumptions C1-C4
and standard Markovian service models, the fork/race/fold framework
recovers selected canonical queueing constructions in the tested cases,
and extends the vocabulary for by adding topology as a control variable.
The executable proofs in §13 include direct tests for Little,
Erlang-style blocking behavior and Jackson-style bottleneck limits

.

\textbf{Precision on ``containment.''} Little's Law () holds under
remarkably weak assumptions --- it requires only ergodicity and finite
expectations, not Markovian arrivals or any specific service-time
distribution (Little \& Graves, 2008

). The C1-C4 conditions are \emph{not} equivalent to Little's Law's
assumptions; they are \emph{stronger} (C3 demands deterministic fold, C4
demands finite termination). Containment means: every system satisfying
C1-C4 with also satisfies the conditions under which Little's Law holds,
and the framework produces the same steady-state predictions. The
converse does not hold --- Little's Law applies to systems that violate
C3 (non-deterministic service) or have no fold semantics at all. What
fork/race/fold adds is not a relaxation of Little's assumptions but an
extension of the \emph{vocabulary}: when , topology becomes a control
variable that queueing theory has no notation for.

Little's Law constrains long-run occupancy/latency averages and is
agnostic to detailed topology. In this manuscript it is exercised on
path-like constructions. When , Little's Law can still hold while
remaining silent about topology-control questions -- how paths interact,
when to fork, when to fold, and when to vent.

The companion suite now adds a \textbf{sample-path conservation
identity} in the finite executable scope: for finite arrival traces and
positive service requirements on a work-conserving single-server queue,
the identity

holds because both sides count the same customer-time in system -- the
left by time slice, the right by job. In the executable model this is
exercised by exhaustively enumerating all finite tick-level
work-conserving service choices on selected small traces (preemption
allowed at tick boundaries), then recovering familiar disciplines such
as FIFO, LIFO, static priority and shortest-remaining-processing-time as
named fold-selection policies inside that larger family. Representative
discretized exponential, Erlang, hyperexponential and lognormal service
families produce the same identity after sampling. In this bounded
scope, queue discipline changes which path folds next, not the
conservation law itself.

The bounded formal layer now extends this to \textbf{finite multi-class
open networks}: jobs belong to distinct classes, classes carry different
routes through the node graph, and finite service-law scenarios vary the
per-stage service realizations. Under node-local work-conserving
dispatch, the same conserved quantity reappears at network scope:

so the invariant is still customer-time in system, now aggregated across
classes and nodes rather than a single queue.

A further bounded stochastic layer treats arrivals, class mixes, route
choices and service realizations as a \textbf{finite-support weighted
scenario family}. Because the customer-time identity is pathwise for
each realization, finite linearity lifts it to expectation in the
executable model:

This remains a finite-support stochastic statement, not a full
probabilistic-process semantics claim.

The next bounded step eliminates that caveat for a tiny state space: an
\textbf{exact finite-state probabilistic transition kernel} evolves the
full probability-mass distribution of a bounded FIFO queue tick by tick,
without pre-expanding the leaves into an external scenario table. In
that kernel the same invariant holds directly at the distribution level,

where is mass-weighted customer-time accumulated through tick , is
mass-weighted departed sojourn, and is mass-weighted open age. This is
an exact finite-state probabilistic semantics result for the bounded
kernel, still short of unbounded or continuous-time queueing theory.

The same move now extends one level outward to a bounded
\textbf{multi-class open-network kernel}: the full probability mass over
class-dependent route states is propagated tick by tick for a two-node
network, and the same distribution-level invariant is rechecked there.
Its longest branch is also a useful witness for the pipeline story's
small-data pathology: a beta job arrives first on node 2, immediately
turns onto node 1, and is followed one tick later by an alpha job that
also starts on node 1. With only two potential arrivals, that
reverse-route collision still stretches completion to six ticks, making
it a minimal executable witness that ramp-up can dominate even the
tiniest nontrivial workload.

That bounded exactness can be pushed one rung higher without leaving the
finite executable regime: a larger \textbf{three-arrival, three-class,
three-node} witness carries the entire arrival cube exactly. The
executable harness evolves the corresponding probability-mass kernel
directly, while the formal layer checks the same weighted conservation
law over the full arrival cube in one stateful model. This does not yet
yield arbitrary exact multiclass/open-network semantics, but it moves
beyond the minimal witness and shows that the exact probabilistic
argument survives a meaningfully larger open-network geometry.

The limit side is now stronger than a schema shell. Constructively,
every finite truncation of a balanced weighted scenario family remains
balanced. Formally, the Lean companion now proves seven lifts or
stationary laws: exact conservation for infinite weighted scenario
families via tsum, direct countably supported stochastic queue laws via
PMF and PMF.toMeasure, exact conservation for measurable queue
observables via lintegral, a monotone truncation-to-limit theorem via
lintegral\_iSup, the stable M/M/1 geometric stationary occupancy law
with finite mean queue length for , a finite-node product-form
open-network occupancy law with exact singleton mass and total mean
occupancy when a stable throughput witness satisfying the Jackson
traffic equations is supplied, and a trajectory-level Cesaro balance
theorem for unbounded open-network sample paths whose residual open age
has a limit. The current in-package route to that witness is now
explicit: start from , form the spectral candidate under , prove it
satisfies the traffic equations and the nodewise bounds , , instantiate
the direct spectral product-form witness, and then use a
Knaster-Tarski-style dominance argument to force the monotone
least-fixed-point witness below the same candidate. The adaptive
extension is now one rung stronger too: the companion mechanizes the
comparison principle from a state-indexed routing family to any
dominating ceiling kernel, including the pointwise supremum kernel when
that ceiling remains substochastic, so adaptive constructive throughput
is trapped below the ceiling's spectral candidate. It also now closes
one concrete raw adaptive family all the way down: a bounded two-node
rerouting model derives its own dominating ceiling kernel, proves the
ceiling is strict-row-substochastic and spectrally contractive, bounds
the adaptive constructive throughput by a closed-form candidate, and
emits a linear drift witness on the bounded adaptive state space. That
same closure now has both an executable mirror and a Lean-exported
witness catalog in the companion package:
\texttt{adaptive-supremum-witness.\{json,md\}} enumerates the bounded
state space and several schedule patterns against the same closed-form
ceiling, while \texttt{formal-adaptive-witness-catalog.\{json,md\}} pins
the concrete raw adaptive witness \texttt{α\ =\ (1/4,\ 11/40)} together
with drift gap \texttt{1/8} and spectral radius \texttt{0}. The
companion also no longer leaves the generic state-dependent frontier
entirely uninstantiated: in addition to the assumption-parameterized
Foster-Lyapunov/irreducibility schema, it now packages bounded vacation,
retrial, and reneging queue families at both the abstract-family and
kernel levels, yielding stationary and terminal balance theorems once
the corresponding drift, irreducibility, and open-age hypotheses are
discharged. This is a genuine measure-theoretic lifting of the
sample-path conservation law into infinite-support or continuous-support
settings, now with constructive Jackson traffic-equation fixed points
and dominance bounds, a bounded Jackson-style product-form layer
grounded in those traffic equations, concrete vacation/retrial/reneging
family instantiations, an adaptive comparison layer for state-indexed
routing families, a concrete raw adaptive ceiling/drift witness,
executable and Lean-exported adaptive witness artifacts, an
assumption-parameterized state-dependent stability interface, and an
ergodic interface for open networks. What it still does not provide is
arbitrary exact multiclass/open-network semantics beyond the current
bounded witnesses, a generic constructive discharge of the adaptive or
state-dependent drift/irreducibility schema for arbitrary raw families,
or a witness-free positive-recurrence derivation for unbounded open
stochastic networks.

The pipeline Reynolds number is used here as a complementary topology
diagnostic (not a replacement for Little's Law):

\textbf{Queueing Theory} \textbar{} \textbf{Fork/Race/Fold} \textbar{}

~(items in system) \textbar{} (parallel paths in system) \textbar{}

Utilization \textbar{} (stages / chunks) \textbar{}

~for stability \textbar{} Heuristic bands in this manuscript:
laminar-like; turbulent-like \textbar{}

M/M/1, M/M/c, M/G/1 variants \textbar{} Laminar, transitional, turbulent
regimes \textbar{}

Arrival rate \textbar{} Fork rate \textbar{}

Service rate \textbar{} Fold rate \textbar{}

Queue discipline (FIFO, priority) \textbar{} Fold strategy (quorum,
weighted, consensus) \textbar{}

In queueing theory, an M/M/1 queue represents the simplest non-trivial
model of a waiting line. It describes a memoryless system with a single
server where arrivals and service times are essentially random. Its
notation follows Kendall's Notation, where each letter defines a
specific characteristic of the system:

\begin{itemize}
\tightlist
\item
  M (Markovian/Memoryless) Arrival: Customers arrive according to a
  Poisson process. This means the time between arrivals follows an
  Exponential distribution. It is ``memoryless'' because the time until
  the next arrival doesn't depend on how much time has already passed.\\
\item
  M (Markovian/Memoryless) Service: The time it takes to serve a
  customer also follows an Exponential distribution.\\
\item
  1 (Single Server): There is only one station or person processing the
  queue.
\end{itemize}

Every M/M/1 queue is a pipeline with , one stage and Poisson arrivals.
The framework does not contradict queueing theory -- it embeds canonical
one-path constructions. When , reduces to utilization. When , adds
topology-aware vocabulary for sequential-to-multiplexed transition,
fork-width tuning and topological mismatch cost.

\subsubsection{\texorpdfstring{\textbf{4.2 Erlang's Formula as Fold
Without
Fork}}{4.2 Erlang's Formula as Fold Without Fork}}\label{erlangs-formula-as-fold-without-fork}

Erlang's B formula gives the blocking probability for servers with no
queue:

In fork/race/fold terms, Erlang's system is a race over servers -- but
without fork. Arrivals are not forked; they are routed to a single
server. The system cannot exploit parallelism because it has no fork
operation. Blocking occurs when all paths are occupied -- but there is
no mechanism to create \emph{new} paths on demand.

While Agner Krarup Erlang provided the mathematical logic that allows us
to build networks that don't collapse under pressure, he didn't have
fork/race/fold.

Fork/race/fold can reduce blocking pressure by making path creation
dynamic. When demand exceeds capacity, fork creates new paths (
increases). When demand subsides, fold and venting remove paths (
decreases). The topology adapts to load. Erlang's formula describes a
\emph{static} case; fork/race/fold models a \emph{dynamic} case.

\subsubsection{\texorpdfstring{\textbf{4.3 Jackson Networks as
Fixed-Topology
Pipelines}}{4.3 Jackson Networks as Fixed-Topology Pipelines}}\label{jackson-networks-as-fixed-topology-pipelines}

James R. Jackson was a mathematician at UCLA who, by 1963, realized
that, in the real world, queues don't exist in isolation: a factory
floor, a hospital, or a data center are all complex networks, not simple
conveyor belts.

Jackson's theorem

proves that open networks of M/M/c queues have product-form stationary
distributions. But Jackson networks have \textbf{fixed topology} -- the
routing matrix is constant. Fork/race/fold has \textbf{dynamic topology}
-- fork creates paths, venting removes them, fold merges them. The
topology is the control variable, not a parameter.

A Jackson network can be represented in this vocabulary as a
fixed-topology case (fixed routing matrix, fixed service structure) with
no dynamic vent policy in the standard formulation. Adding dynamic
routing, load-dependent forking, or failure-driven path removal moves
beyond classical Jackson assumptions.

You enter the domain of fork/race/fold, where topology is treated as a
variable rather than a fixed parameter.

\subsubsection{\texorpdfstring{\textbf{4.4 What Replaces
What}}{4.4 What Replaces What}}\label{what-replaces-what}

Queueing theory asks: \emph{given a fixed topology, what is the
steady-state behavior?}

Fork/race/fold asks: \emph{what topology should the system have at each
decision point?}

The Reynolds number provides a runtime heuristic for this question. In
the benchmarked regime bands used here: suggests sequential sufficiency,
suggests multiplexing opportunity, and suggests widening fork degree.
The topology is not fixed; it is adapted from the same measurement that
drives scheduling.

This contrast is used as a heuristic: queueing theory emphasizes
steady-state behavior for fixed topologies, while fork/race/fold
emphasizes topology-adaptation decisions under the assumptions used
here.

\subsection{\texorpdfstring{\textbf{5. The Quantum Vocabulary Is
Structural}}{5. The Quantum Vocabulary Is Structural}}\label{the-quantum-vocabulary-is-structural}

The following correspondences are structural mappings between
quantum-mechanical operations and computational operations, with
photosynthetic antenna complexes (§1.5) as a literal quantum anchor
example. In §6.12, I show that the Feynman path integral admits a
fork/race/fold interpretation within this abstraction.

\textbf{Relation to prior formalisms.} The concurrent-computation
literature offers several models with overlapping expressiveness. Petri
nets (Petri, 1962

) represent fork as transition firing and fold as place merging; they
excel at deadlock analysis but lack native race and vent semantics. The
-calculus (Milnor, 1999

) models dynamic channel creation (akin to fork) and synchronization
(akin to fold) with full compositionality; it does not, however, expose
a topological characterization (, covering spaces) or a thermodynamic
framing. Speculative execution in CPU microarchitectures (Tomasulo, 1967

; Smith \& Sohi, 1995

) implements fork (issue multiple paths), race (retire the correct path
first), and vent (flush mispredicted paths) at the hardware level ---
the closest engineering analogue to fork/race/fold, discovered
independently by processor designers optimizing instruction-level
parallelism. Byzantine fault-tolerant consensus protocols (Castro \&
Liskov, 1999

; Yin et al., 2019

) implement quorum fold under adversarial conditions, with explicit vent
of Byzantine-faulty replicas. Fork/race/fold does not replace these
formalisms. It provides a unifying topological and thermodynamic
vocabulary --- , , the First Law --- that situates them within a common
framework and highlights structural properties (the deficit diagnostic,
the Reynolds-number phase transition) that are not jointly explicit in
any single prior formalism here.

\textbf{Quantum Operation} \textbar{} \textbf{Computational Operation}
\textbar{} \textbf{What It Does} \textbar{}

\textbf{Superposition} \textbar{} Fork \textbar{} paths exist
simultaneously, outcome undetermined \textbar{}

\textbf{Measurement} \textbar{} Observe \textbar{} Non-destructive state
inspection without triggering fold \textbar{}

\textbf{Collapse} (QM term) \textbar{} Race / Fold \textbar{} Resolve to
a definite state \textbar{}

\textbf{Tunneling} \textbar{} Early exit \textbar{} Bypass remaining
computation when a path is conclusive \textbar{}

\textbf{Interference} \textbar{} Consensus \textbar{} Constructive:
agreeing signals amplify. Destructive: disagreeing signals cancel
\textbar{}

\textbf{Entanglement} \textbar{} Shared state \textbar{} Correlated
streams that see each other's mutations \textbar{}

\subsubsection{\texorpdfstring{\textbf{5.1
Superposition}}{5.1 Superposition}}\label{superposition}

After fork, a computation exists in simultaneous states -- the outcome
is undetermined until fold. This is computational superposition. It has
a closely related structural form to quantum superposition: a quantum
state is a superposition of basis states, and a forked computation is a
superposition of branch states. Fold then projects to a definite outcome
in the computational model.

In photosynthetic antenna complexes (§1.5), this is literal quantum
superposition. The excitation energy is in a quantum superposition
across multiple pigment molecules until it collapses at the reaction
center. The fork() operation serves as the computational analogue in
this framework.

\subsubsection{\texorpdfstring{\textbf{5.2
Tunneling}}{5.2 Tunneling}}\label{tunneling}

In quantum mechanics, tunneling allows a particle to pass through a
potential barrier that classical physics says is impassable. In
fork/race/fold, tunneling allows a computation to bypass the ``barrier''
of waiting for all paths to complete.

A tunnel predicate fires when a single path's result is conclusive
enough that remaining paths are irrelevant. It's worth reiterating here
again that this is different from race (which picks the \emph{fastest})
and different from fold (which waits for \emph{all}). Tunneling picks
the \emph{first sufficient result} and vents everything else -- it
``tunnels through'' the waiting barrier.

Tunneling is not a fifth primitive. It is a composition: race(predicate)
+ vent(losers) --- race with a quality predicate instead of a speed
predicate. Topologically, tunneling operates on homotopy-equivalent
paths (§3.2) but selects by a criterion other than arrival time. Where
race exploits temporal homotopy (all paths reach the same destination,
pick the fastest), tunneling exploits quality homotopy (all paths
produce valid results, pick the first that's sufficient). The fallback
to race or fold when the predicate is too strict confirms this:
tunneling degrades gracefully into its constituent primitives.

Use case: a diagnostic pipeline forks blood test, MRI and genetic
screening. The blood test returns a conclusive positive. Tunneling
fires: the MRI and genetic screening are vented. No need to wait. The
tunnel predicate evaluated quality, not speed.

\subsubsection{\texorpdfstring{\textbf{5.3
Interference}}{5.3 Interference}}\label{interference}

Constructive interference amplifies signals that agree. Destructive
interference cancels signals that disagree. In fork/race/fold, the
consensus fold strategy implements both:

\begin{itemize}
\tightlist
\item
  \textbf{Constructive}: compute pairwise agreement across all results.
  Values where streams agree are amplified (kept). This is signal
  extraction from noise.\\
\item
  \textbf{Destructive}: values where streams agree are kept. This is
  outlier detection -- finding the signal that \emph{disagrees} with the
  majority.
\end{itemize}

The pairwise comparison is the interference pattern. The resulting fold
is the detected signal.

\subsubsection{\texorpdfstring{\textbf{5.4
Entanglement}}{5.4 Entanglement}}\label{entanglement}

In quantum mechanics, entangled particles share state across arbitrary
distances -- measuring one instantly determines the other even without
shared communication. In fork/race/fold, entangled streams share a
mutable reference. Mutations by one stream are visible to all others. No
locks, no synchronization -- the shared state is the entanglement.

Use case: vote tallying. Fork streams to count ballot boxes. All streams
share an accumulator. Each stream's partial count is immediately visible
to monitoring (measurement) without triggering fold.

\subsubsection{\texorpdfstring{\textbf{5.5
Measurement}}{5.5 Measurement}}\label{measurement}

Measurement in quantum mechanics famously disturbs the system --
measuring folds the superposition. In fork/race/fold, measurement is
\textbf{non-destructive}: you can observe the current state of all
forked streams without triggering fold or venting. The distinction is
intentional -- I want observability without interference.

The measure() operation returns a snapshot: stream states, intermediate
results, timing. Monitoring dashboards, progress bars, debugging -- all
measurement, all non-destructive.

\subsection{\texorpdfstring{\textbf{6. The Geometry of
Failure}}{6. The Geometry of Failure}}\label{the-geometry-of-failure-1}

The topology (§3) classifies the \emph{shape} of computation. The
queueing containment (§4) situates it within existing theory. The
quantum vocabulary (§5) names its operations. This section provides the
stability geometry: fork/race/fold operates as a manifold of
\textbf{curvature, drift, sinks, and voids}. In this framing, failure is
not an accident. It is a state transition governed by a conservation
law.

\subsubsection{\texorpdfstring{\textbf{6.1 The Geometric
Dictionary}}{6.1 The Geometric Dictionary}}\label{the-geometric-dictionary}

\textbf{Primitive} \textbar{} \textbf{Geometric Analogue} \textbar{}
\textbf{Symbol} \textbar{}

\textbf{Fork} \textbar{} Curvature injection / stress widening
\textbar{} \textbar{}

\textbf{Race} \textbar{} Drift along the manifold \textbar{} \textbar{}

\textbf{Fold} \textbar{} Projection to the stable sink \textbar{}
\textbar{}

\textbf{Vent} \textbar{} Fail-flow boundary map / void ledger \textbar{}
\textbar{}

\textbf{Backpressure} \textbar{} Restorative gravity \textbar{}
\textbar{}

\textbf{Stream} \textbar{} Geodesic through the cover space \textbar{}
\textbar{}

\textbf{Frame} \textbar{} Local coordinate chart \textbar{} \textbar{}

The governing question is now:

If outside a bounded small set, the system drifts back toward stability.
If , the manifold no longer bends inward strongly enough to recover.

\subsubsection{\texorpdfstring{\textbf{6.2 Fork as Curvature
Injection}}{6.2 Fork as Curvature Injection}}\label{fork-as-curvature-injection}

A fork creates parallel paths. Geometrically, it widens the manifold and
increases the exposed failure surface. Each new path adds optionality,
but it also adds stress because more of the state space must later be
reconciled.

This is why matters. Each independent cycle counted by is a unit of
exposed topological freedom, and each unclosed degree of freedom
contributes to structural brittleness. The \textbf{Bule} functions as a
\textbf{stress unit}:

A system at 10 B is not merely ``10 paths behind.'' It is carrying 10
units of unmodeled curvature: its flow cannot see the failure surface it
must later traverse.

\subsubsection{\texorpdfstring{\textbf{6.3 Race as
Drift}}{6.3 Race as Drift}}\label{race-as-drift}

Race is the manifold in motion. Each active path advances, stalls, or
degrades the system's return toward the origin. The geometric question
is not which path is fastest in isolation but whether the
\textbf{expected next step} of the whole transition kernel still points
inward.

This is the role of the Foster-Lyapunov drift schema. Let measure
distance from the stable sink for state -- for example, unresolved
frontier, outstanding fold debt, and exposed . Then:

Negative drift means restorative pressure. Positive drift means failure
expansion. In queue-like regimes the margin can be written:

where is arrival pressure, is service pressure, and is restorative
shedding pressure. When , failure is not ``bad luck.'' The field itself
now points away from the safe region.

\subsubsection{\texorpdfstring{\textbf{6.4 Fold as Return to the
Sink}}{6.4 Fold as Return to the Sink}}\label{fold-as-return-to-the-sink}

Fold is the moment the system pays its reconciliation debt and projects
back toward a stable sink. In a single-survivor setting, this projection
is necessarily lossy: many live branches become one canonical state.

That gives the bounded theorem schema I refer to here as
\textbf{THM-FAIL-MINCOST}:

for any finite -branch computation that collapses to one surviving state
in finite time. The ledger can be paid as vented branches, merged debt,
deferred cleanup, or explicit reconciliation work, but it cannot be paid
as zero. Collapse has a floor.

This reframes the binary question ``did the race succeed?'' into a
geometric one: \textbf{how did the manifold pay to become single-valued
again?}

\subsubsection{\texorpdfstring{\textbf{6.5 Venting as
Fail-Flow}}{6.5 Venting as Fail-Flow}}\label{venting-as-fail-flow}

Venting is not merely waste. It is the boundary operator that keeps the
Void bounded. When a path cannot contribute useful stability, vent maps
that path into the failure ledger and prevents it from dragging siblings
outward.

This is the operational meaning of \textbf{Failure-Modulated Flow
Validation (FMFV)}: the most reliable systems are the ones that let
failure reshape the flow before the whole manifold escapes. A reneging
customer, a canceled speculative branch, or a discarded codec result is
not a defect in this framing. It is restorative geometry. The system
fails \emph{some work} so the system does not fail \emph{as a system}.

Fail-safe systems try to forbid failure. \textbf{Fail-flow systems} pay
failure locally so the remaining graph keeps negative drift globally.

\subsubsection{\texorpdfstring{\textbf{6.6 Backpressure as Restorative
Gravity}}{6.6 Backpressure as Restorative Gravity}}\label{backpressure-as-restorative-gravity}

Backpressure is the mechanism by which the manifold regains inward
curvature. When input flow exceeds processing capacity, the state moves
uphill away from the sink. Backpressure changes the field by lowering ,
increasing effective , or both.

This is why stability can be described as a potential well. The small
set is the basin of attraction. Recovery is gravity. Failure is the
moment the combined restorative gravity is no longer large enough to
overcome arrival pressure .

In the rotational frame (the Worthington Whip), this same effect appears
as tension release: widening the manifold creates debt, and the whip
snap pays it in one controlled inward motion.

\subsubsection{\texorpdfstring{\textbf{6.7 The Failure
Trilemma}}{6.7 The Failure Trilemma}}\label{the-failure-trilemma}

The geometry yields a hard trilemma for finite single-survivor systems:

\begin{enumerate}
\def\labelenumi{\arabic{enumi}.}
\tightlist
\item
  Collapse to one surviving state.\\
\item
  Pay zero loss, zero debt, and zero discarded alternatives.\\
\item
  Converge in finite time.
\end{enumerate}

You may choose at most two.

If you demand finite-time collapse to one state, the Void must record a
payment. If you refuse all payment, the system stays forked forever.
Reliability is therefore not the elimination of failure. Reliability is
the discipline of \textbf{choosing how failure is paid}.

\subsubsection{\texorpdfstring{\textbf{6.8 The Information-Theoretic
Void}}{6.8 The Information-Theoretic Void}}\label{the-information-theoretic-void}

\begin{description}
\tightlist
\item[The Shannon connection acts as a map of the Void]
\end{description}

\begin{itemize}
\tightlist
\item
  \textbf{Fork} creates up to bits of alternative futures under
  uniform-path assumptions.\\
\item
  \textbf{Race} reveals which futures remain viable.\\
\item
  \textbf{Fold} compresses many futures into one retained outcome.\\
\item
  \textbf{Vent} records the discarded alternatives.
\end{itemize}

The ledger identity is:

The Void is therefore not the unknown. It is the accounting surface for
branches that existed and were then sacrificed so the system could
re-enter a stable sheet.

\subsubsection{\texorpdfstring{\textbf{6.9 The Pipeline as a Stability
Surface}}{6.9 The Pipeline as a Stability Surface}}\label{the-pipeline-as-a-stability-surface}

The Triangle (§0.1) is a stability surface:

\begin{itemize}
\tightlist
\item
  \textbf{Ramp-up (fork):} curvature increases as work enters the
  manifold.\\
\item
  \textbf{Plateau (race):} drift is balanced and the system occupies its
  widest stable sheet.\\
\item
  \textbf{Ramp-down (fold):} the manifold contracts back toward the
  sink.
\end{itemize}

Idle slots are holes in that surface. Turbulent multiplexing (§7.2)
fills those holes by restoring inward pressure with other work. The
Worthington Whip (§7.3) redistributes curvature so the system pays its
reconciliation debt in a sharp but controlled return rather than as
diffuse instability.

\subsubsection{\texorpdfstring{\textbf{6.10 Three Geometric
Laws}}{6.10 Three Geometric Laws}}\label{three-geometric-laws}

\textbf{First Law (failure conservation).} Failure cost does not
disappear. It is redistributed between retained work and the Void
ledger.

\textbf{Second Law (negative drift or instability).} Stability requires
outside a bounded small set. Without inward drift, the manifold escapes.

\textbf{Third Law (bounded Void).} A reliable system is one whose vent
map bounds the Void. If the boundary fails, the failure surface expands
without limit.

The complete mapping:

\textbf{Fork/Race/Fold} \textbar{} \textbf{Geometric Mechanics}
\textbar{} \textbf{Stability Law} \textbar{}

Fork \textbar{} Curvature injection \textbar{} Widens the exposed
failure surface \textbar{}

Race \textbar{} Drift \textbar{} Determines inward vs outward motion
\textbar{}

Fold \textbar{} Sink projection \textbar{} Pays reconciliation debt to
re-enter stability \textbar{}

Vent \textbar{} Void boundary map \textbar{} Keeps local failure from
globalizing \textbar{}

~\textbar{} Parallel-path curvature \textbar{} Measures how much
structure must be reconciled \textbar{}

~\textbar{} Topological stress \textbar{} Quantifies unseen failure
surface \textbar{}

~\textbar{} Discarded alternatives ledger \textbar{} Records what
collapse could not retain \textbar{}

~\textbar{} Restorative gravity \textbar{} Positive margin implies
recoverability \textbar{}

Pipeline Triangle \textbar{} Stability surface \textbar{} Shape of
exposure and return \textbar{}

\subsubsection{\texorpdfstring{\textbf{6.11 Transformers Are
Fork/Race/Fold
Graphs}}{6.11 Transformers Are Fork/Race/Fold Graphs}}\label{transformers-are-forkracefold-graphs}

The geometric failure framing highlights that convolutional neural
networks and transformers can be represented as fork/race/fold graphs at
useful levels of abstraction.

\textbf{Multi-head attention is fork/race/fold.} The input splits into
heads (each with its own , , projections). This is fork: . All heads
compute attention over the same sequence simultaneously -- race.
Concatenation plus linear projection -- fold: the merger function that
produces a single representation. Softmax suppression (low-attention
scores ) is continuous venting: the system shedding paths that don't
contribute.

\textbf{Feed-forward layers are fork/fold.} The input expands from to --
fork into a 4x wider representation. The activation function (ReLU,
GELU) performs \emph{soft venting}: zeroing or suppressing
non-contributing neurons. The contraction back to is fold. The
distinction from computational vent is important: in fork/race/fold,
vent is \emph{irreversible} --- a vented path is structurally removed
and cannot contribute to fold. In FFN layers, ReLU-zeroed neurons are
structurally present (their weights persist and gradient descent can
reactivate them in subsequent forward passes). The FFN vent is therefore
\emph{per-inference} irreversible but \emph{per-training-step}
reversible --- a softer form of the primitive. During inference (the
thermodynamic ``measurement''), the zeroed activations are genuinely
vented: they contribute zero to the fold and their potential energy is
dissipated. During training, the vent boundary shifts as gradients
update the weights that determine which neurons fire.

\textbf{Residual connections are fork with two paths.} The skip
connection and the transformed path: fork(identity, transform). Addition
is fold via sum. .

\textbf{CNNs follow the same pattern per spatial region.} filters
applied to the same receptive field is fork. All filters compute
simultaneously is race. Pooling is fold -- and max pooling is literally
winner-take-all fold.

\textbf{Transformer Component} \textbar{} \textbf{Primitive} \textbar{}
\textbf{β1\hspace{0pt}} \textbar{} \textbf{Stability Role} \textbar{}

Multi-head attention ( heads) \textbar{} fork/race/fold \textbar{}
\textbar{} Parallel paths whose scores determine inward drift \textbar{}

FFN expansion () \textbar{} fork/vent/fold \textbar{} 3 \textbar{}
Widen, prune, and reproject to a stable representation \textbar{}

Residual connection \textbar{} fork/fold \textbar{} 1 \textbar{}
Two-path fork with additive reconciliation \textbar{}

Softmax attention \textbar{} continuous vent \textbar{} -- \textbar{}
Gradually suppress unstable contributions \textbar{}

Dropout \textbar{} stochastic vent \textbar{} -- \textbar{} Random path
removal during training \textbar{}

Layer norm \textbar{} measure \textbar{} 0 \textbar{} Non-destructive
observation of statistics \textbar{}

MoE routing ( of experts) \textbar{} fork/race/vent/fold \textbar{}
\textbar{} Vent unused experts to keep the route bounded \textbar{}

The entire transformer is a \textbf{nested} fork/race/fold/vent graph.
Each layer is fork/fold. Each attention computation within a layer is
fork/race/fold. The stack of layers is a pipeline.

Transformer architecture can be interpreted as a recursive
Wallington-style composition.

\textbf{Backpropagation as drift correction (interpretive lens).} The
loss function can be read as a stability proxy: how far the network's
current fold rule is from the desired sink. The gradient indicates how
to change parameters so future passes recover more reliably and vent
less destructively. Training then appears as iterative repair of the
model's drift field.

\textbf{Mixture of Experts makes the topology explicit.} MoE routing
with experts, top- selection: fork to experts (), race the router's
gating scores, fold the top- results, vent the remaining . The router
\emph{is} the race primitive. The gating function \emph{is} the fold
function. The unused experts \emph{are} vented paths. The sparse
activation pattern \emph{is} the vent ratio . What the ML community
calls ``conditional computation'' is what this paper calls
fork/race/fold with selective venting.

\subsubsection{\texorpdfstring{\textbf{6.12 Structural Correspondences
with Physical
Formalisms}}{6.12 Structural Correspondences with Physical Formalisms}}\label{structural-correspondences-with-physical-formalisms}

The geometric framing is used as a cross-domain mapping to physics. Two
results from fundamental physics are used as structural correspondences
with fork/race/fold, with limited quantitative anchors in cited scope.

\paragraph{\texorpdfstring{\textbf{The Feynman Path Integral (Grade
B+)}}{The Feynman Path Integral (Grade B+)}}\label{the-feynman-path-integral-grade-b}

In quantum electrodynamics, the probability amplitude for a particle
traveling from point to point is:

where is the action along each path. The path-integral calculation is
interpreted here in four phases:

\begin{enumerate}
\def\labelenumi{\arabic{enumi}.}
\tightlist
\item
  \textbf{Fork.} The particle enters all possible trajectories
  simultaneously. Each trajectory is a path with phase . This is fork:
  one input innumerable paths. .\\
\item
  \textbf{Race.} Each path propagates with its own phase accumulation.
  No path ``knows'' about the others during propagation (allowing for
  transport gains). This is race: parallel, independent, timeless
  (unitary evolution is time-reversible).\\
\item
  \textbf{Fold.} The amplitudes sum. Constructive interference
  concentrates amplitude on the classical path (stationary phase). This
  is fold: many paths one probability amplitude. .\\
\item
  \textbf{Vent.} Destructive interference eliminates non-classical
  paths. Their amplitudes cancel to zero. This is vent: paths that
  contribute no useful work are dissipated. ``Propagate down, never
  across'' -- destructively interfered paths do not affect the surviving
  amplitude.
\end{enumerate}

The classical limit () recovers the path of stationary action -- the
unique classical trajectory. This is the subspace: one path, no fork, no
race, no vent. In this abstraction, classical mechanics is represented
as the limit of the mapping, just as sequential pipelines are the
degenerate case of the Wallington Rotation.

\textbf{This is a structural mapping with explicit boundaries.} The path
integral can be mapped to a fork/race/fold interpretation: the sum over
paths maps to fork, interference maps to fold/vent, and the stationary
phase approximation maps to the projection. Feynman diagrams are
computation graphs whose topological properties ( = loop order) track
calculation difficulty, similar to how tracks pipeline complexity in §3.

\textbf{Validated boundary condition.} The correspondence is
operationally exact only in the linear full-aggregation regime. Five
companion validations make that boundary explicit. In the finite-kernel
unit harness
(companion-tests/src/quantum-correspondence-boundary.test.ts)

, linear fold reproduces discrete path-sum evolution exactly (kernel
composition equals explicit path enumeration), preserves partition
additivity, and remains permutation-invariant on the cancellation
witness; winner-take-all and early-stop folds fail those same checks. In
the fold-ablation harness
(companion-tests/src/quantum-recombination-ablation.test.ts) and its
reproducible artifact writer
(companion-tests/scripts/quantum-recombination-ablation.ts, output
companion-tests/artifacts/quantum-recombination-ablation.\{json,md\})

, the path family is held fixed while only the recombination rule is
swapped: the predicted loss matrix is recovered exactly, with linear
fold preserving kernel agreement, partition additivity, order invariance
and cancellation, while winner-take-all and early-stop each show
kernel-agreement distance 0.354, partition/order distance 2.000, and
cancellation magnitude\(^2\) 1.000. In the Lean theorem package
(companion-tests/formal/lean/Lean/ForkRaceFoldTheorems/Claims.lean)

, the algebraic skeleton of the same boundary is mechanized in a minimal
integer-valued model: linear fold is globally partition-additive,
preserves the cancellation target family x + (-x) = 0, and is equivalent
to the more general cancellation-difference family fold(x,-y) = x - y;
any non-additive fold must miss some member of that family, and
winner-selection/early-stop miss the concrete x + (-x) witness. A
Lean-emitted witness catalog
(companion-tests/artifacts/formal-witness-catalog.\{json,md\}) now
exports 7 concrete cancellation, partition, and order counterexamples,
and quantum-correspondence-boundary.test.ts consumes those exported
witnesses directly rather than hardcoding them. A seeded Gnosis
cancellation benchmark
(companion-tests/artifacts/gnosis-fold-training-benchmark.\{json,md\})
then keeps topology, parameter count, and data fixed across three .gg
programs and changes only the fold strategy: linear fold reaches eval
MSE 0.000 with 95\% seed-bootstrap interval {[}0.000, 0.000{]}, while
winner-take-all and early-stop settle at 0.408/0.735 with intervals
{[}0.396, 0.421{]}/{[}0.732, 0.740{]}; cancellation-line absolute error
is 0.000, 0.834, and 0.764. A paired seeded negative-control benchmark
(companion-tests/artifacts/gnosis-negative-controls.\{json,md\}) then
keeps the same topologies but moves to one-path target families where no
cross-path cancellation or dual-expert summation is required; there the
separation disappears exactly as predicted, with affine-left-only and
positive-x single-expert controls both yielding max inter-strategy
eval-MSE gap 0.000 and min exact-within-tolerance rate 1.000. Finally, a
harder seeded Gnosis mini-MoE routing benchmark
(companion-tests/artifacts/gnosis-moe-routing-benchmark.\{json,md\})
keeps a four-expert routed topology and fixed 16-parameter budget while
swapping only the fold strategy: linear fold reaches eval MSE 0.001,
winner-take-all 0.328, and early-stop 0.449, with 95\% seed-bootstrap
intervals {[}0.001, 0.001{]}, {[}0.267, 0.389{]}, and {[}0.444,
0.457{]}; the dual-active-region absolute error is 0.027, 0.402, and
0.474. The assembled manuscript figures are emitted automatically to
companion-tests/artifacts/ch17-correspondence-boundary-figure.\{json,md,svg\}
and
companion-tests/artifacts/ch17-boundary-expansion-figure.\{json,md,svg\},
and the full evidence bundle is fingerprinted in
companion-tests/artifacts/ch17-replication-pack.\{json,md\} with the
one-command outside rerun surface bun run
test:ch17-external-replication, 34 manifest entries, and 14 generated
artifacts. So the shared structure is ``fork, independent propagation,
recombination to one output,'' but the recombination mechanics differ:
physical path integrals sum linearly; computational winner/race folds
select nonlinearly.

\paragraph{\texorpdfstring{\textbf{The Physics Hierarchy: Progressive
Folds}}{The Physics Hierarchy: Progressive Folds}}\label{the-physics-hierarchy-progressive-folds}

In this abstraction, the path integral, the Schrödinger equation, and
Newton's laws can be arranged as an interpretive hierarchy of
progressively coarser folds. This is a modeling view, not a claim of
full formal equivalence.

\textbf{Level 0: The Path Integral (full fork/race/fold).} All paths.
All interferences. No approximation. .

\textbf{Level 1: The Schrödinger Equation (race-like differential
dynamics).} Feynman showed

that evaluating the path integral in the limit of infinitesimal time
steps recovers the Schrödinger equation in the standard derivation:

In this framing, this is what happens when the path-integral evolution
is expressed as a \emph{local} differential equation in the standard
derivation. The wave function is the bookkeeping device that tracks the
superposition of all racing paths at each instant. is the probability
density -- the energy distribution across surviving paths.

In physics and mathematics, the Hamiltonian is a mathematical operator
that represents the total energy of a system. It's a function that sums
up all the energy ``bank accounts'' of a particle or system. In this
mapping, the Hamiltonian is interpreted as the race operator: it governs
how potential converts to kinetic at each infinitesimal step.

In this mapping, the Schrödinger equation is a race-like local dynamics
equation. It is treated as a local form of global path exploration in
this abstraction. The wave function carries information about which
paths are still active and with what amplitude.

\textbf{Quantized energy levels are fold constraints.} For bound systems
(electrons in atoms, particles in wells), the Schrödinger equation
admits only discrete solutions -- specific energy eigenvalues. These are
not inputs to the equation; they \emph{emerge} from the fold boundary
conditions. The requirement that at infinity (the wave function must be
normalizable) is a fold constraint: it eliminates all solutions that
don't converge. The surviving eigenvalues are the fold results. Lasers,
LEDs, atomic clocks and MRI machines all depend on these quantized fold
outputs.

\textbf{Quantum tunneling as incomplete venting.} Classically, a
particle encountering a potential barrier higher than its kinetic energy
is disallowed from crossing. But the Schrödinger equation shows that
decays exponentially through the barrier rather than dropping to zero.
If the barrier is thin enough, nonzero amplitude leaks through. In this
framework, that is an incomplete-vent analogue. Flash memory, scanning
tunneling microscopes, and nuclear fusion in stars exploit this effect.

\textbf{Level 2: Stationary Phase Approximation (the vent operator).} In
the classical limit (), the phase oscillates infinitely fast. Nearly all
paths cancel by destructive interference -- they are vented. Only paths
near the stationary point of the action survive:

The stationary phase approximation acts like a maximal vent operator in
this mapping. It destroys most path information except near-classical
trajectories. . The void () grows as many quantum paths are canceled.

\textbf{Level 3: Newton's Laws (, fully folded).} One path.
Deterministic. No fork, no race, no vent. is the maximally folded result
of the path integral. Classical mechanics is not ``wrong'' -- it is the
degenerate case, just as sequential pipelines are the degenerate case of
the Wallington Rotation.

The hierarchy:

\textbf{Level} \textbar{} \textbf{Theory} \textbar{}
\textbf{Fork/Race/Fold Role} \textbar{} \textbf{β1\hspace{0pt}}
\textbar{} \textbf{Information} \textbar{}

0 \textbar{} Path integral \textbar{} Full engine \textbar{} \textbar{}
All paths, all phases \textbar{}

1 \textbar{} Schrödinger equation \textbar{} Differential race
\textbar{} Finite \textbar{} Wave function tracks superposition
\textbar{}

2 \textbar{} Stationary phase \textbar{} Maximal vent \textbar{}
\textbar{} Only near-classical paths survive \textbar{}

3 \textbar{} Newton's laws \textbar{} Fully folded \textbar{} \textbar{}
One path, deterministic \textbar{}

Each level can be read as a fold/coarse-graining step. Each step
discards information and increases abstraction. In that sense, the
classical tower can be interpreted as nested fold operations on
path-integral structure. Reconstructing finer levels requires
reintroducing information.

This mirrors an information-discard perspective analogous to
coarse-graining under the second-law lens used in this manuscript.

\textbf{Band theory is the covering space formalism applied to
crystals.} When the Schrödinger equation is solved for electrons in a
periodic lattice (silicon, germanium), Bloch's theorem states that
solutions have the form where has the periodicity of the lattice. The
periodic lattice is the base space. The electron's wave function in the
full crystal is the covering space. Bloch's theorem is the covering map
(§3.3) -- it relates the global behavior (energy bands) to the local
structure (unit cell). The band gap -- the energy range where no
electron states exist -- is the void (). Modern semiconductors,
transistors and solar cells rely on this structure.

\paragraph{\texorpdfstring{\textbf{The Virial Theorem (Grade
A-)}}{The Virial Theorem (Grade A-)}}\label{the-virial-theorem-grade-a-}

For self-gravitating systems in equilibrium (gas clouds, galaxies, star
clusters), the virial theorem states:

Half the gravitational potential energy becomes kinetic energy (thermal
motion, radiation). The virial theorem is a constraint on the
\emph{equilibrium state}, not a description of the process that reaches
it. But the process of reaching equilibrium --- gravitational collapse
--- \emph{is} fork/race/fold, and the virial theorem constrains the
energy partition of the fold result. A collapsing gas cloud:

\begin{enumerate}
\def\labelenumi{\arabic{enumi}.}
\tightlist
\item
  \textbf{Fork.} Gravitational potential energy is stored in the spatial
  distribution of mass. Every particle has a trajectory it \emph{could}
  follow. .\\
\item
  \textbf{Race.} Free-fall collapse. Particles accelerate toward the
  center. conversion.\\
\item
  \textbf{Fold.} A star forms -- the bound state. Useful work is
  extracted as nuclear fusion becomes possible. Hydrostatic equilibrium
  is the fold: gravitational compression balanced by radiation
  pressure.\\
\item
  \textbf{Vent.} In this virial-budget interpretation, an order-half
  energy partition appears as dissipative output during relaxation. The
  Kelvin-Helmholtz mechanism is the vent analogue.
\end{enumerate}

The virial theorem gives the exact split: , , therefore . This is a
specific, testable prediction that the First Law produces when applied
to gravity: with the virial theorem constraining the partition ratio.

Fork/race/collapse provides an interpretive description of star
formation that is aligned with measurable physical outcomes.

\paragraph{\texorpdfstring{\textbf{The Weak Force as Vent Operator
(Grade
B+)}}{The Weak Force as Vent Operator (Grade B+)}}\label{the-weak-force-as-vent-operator-grade-b}

Beta decay: . The neutrino carries away energy that is effectively not
recovered locally because it weakly interacts and propagates away. This
is a venting analogue: unstable nuclear configurations dissipate excess
energy toward more stable states.

Supernovae are the extreme case: 99 percent of the gravitational binding
energy ( J) is carried away by neutrinos. The visible explosion --
light, shock wave, ejecta -- is only percent. The vent-to-work ratio: .
Thermodynamic efficiency . In this mapping, the weak interaction acts as
a strong vent analogue.

\paragraph{\texorpdfstring{\textbf{Color Confinement as Anti-Vent (Grade
B)}}{Color Confinement as Anti-Vent (Grade B)}}\label{color-confinement-as-anti-vent-grade-b}

The strong force exhibits a property with no close analogue in the other
nine connections. If you try to separate two quarks (attempt to vent a
color-charged path), the energy stored in the color field creates new
quark-antiquark pairs. Attempted vent automatic fork. In this mapping,
the strong force behaves like an anti-vent via forced forking.

In particle physics, Color Confinement is the phenomenon under which
isolated color-charged quarks or gluons are not observed. Quarks are
locked inside composite particles like protons and neutrons. To be
clear, the ``Color'' in the name refers to Color Charge, which has
nothing to do with visual light; it is the strong force equivalent of
electric charge.

Color confinement can be interpreted as a topological-closure
constraint: isolated color-charged states are not observed, and
attempted separation drives pair creation. In this framework, that
appears as anti-vent behavior.

This strengthens the mapping intuition: the strong-force case behaves
like an anti-vent operator under this vocabulary, while computation
typically permits explicit venting.

\paragraph{\texorpdfstring{\textbf{Symmetry Breaking as Fold (Grade
B+)}}{Symmetry Breaking as Fold (Grade B+)}}\label{symmetry-breaking-as-fold-grade-b}

The Higgs mechanism: above the electroweak energy scale ( GeV), the
electromagnetic and weak forces are unified. Below it, the Higgs field
selects one vacuum state from a continuous family of equivalent states.
The Mexican hat potential is a fork/race/fold landscape:

\begin{itemize}
\tightlist
\item
  \textbf{Fork:} The symmetric state at the top of the potential (all
  vacuum directions equivalent)\\
\item
  \textbf{Race:} The field rolls down the brim (explores vacuum
  states)\\
\item
  \textbf{Fold:} Settles into one minimum (symmetry broken, particles
  acquire mass)\\
\item
  \textbf{Vent:} Goldstone bosons carry away the broken symmetry degrees
  of freedom (three of four are ``eaten'' by the and bosons, becoming
  their longitudinal polarization)
\end{itemize}

Spontaneous symmetry breaking is fold: many equivalent states one
selected state. The void () is the set of unchosen vacua. In this
analogy, observed particle masses correspond to fold-selected outcomes
after symmetry breaking.

\paragraph{\texorpdfstring{\textbf{The Arrow of Time as Fork/Fold
Asymmetry (Grade
B)}}{The Arrow of Time as Fork/Fold Asymmetry (Grade B)}}\label{the-arrow-of-time-as-forkfold-asymmetry-grade-b}

The second law of thermodynamics --- entropy increases over time --- can
be related to fork/fold asymmetry. Fork is reversible in principle if
immediately recombined. Fold is effectively irreversible in this model:
once a winner is selected and losers are vented, discarded path
information is unavailable. The irreversibility enters at the fold/vent
boundary --- the moment of selection. In this interpretation, time's
arrow aligns with movement from (many paths) toward (selected outcome).

\paragraph{\texorpdfstring{\textbf{The Computational Domain as Fold
(Grade
B+)}}{The Computational Domain as Fold (Grade B+)}}\label{the-computational-domain-as-fold-grade-b}

The computational domain can be viewed as a fold boundary that
constrains reachable states and enforces closure in the modeled graph.

\subsubsection{\texorpdfstring{\textbf{6.13 The Optimality
Diagnostic}}{6.13 The Optimality Diagnostic}}\label{the-optimality-diagnostic}

If fork/race/fold is a recurrent shape in finite systems that satisfy
this paper's conservation, irreversibility and minimum-overhead
assumptions, then finding this shape is evidence consistent with
near-optimal topological fit under those assumptions. Not finding it --
where the problem's intrinsic topology demands it -- is a diagnostic for
waste.

Measuring waste in computational systems requires specifying a modeled
structure for both the problem and the implementation. The topological
deficit is the difference between the modeled intrinsic Betti number and
the actual Betti number. This deficit represents potentially unexploited
parallelism in that model.

This opportunity has seen less emphasis because the field has
traditionally focused on algorithmic complexity rather than topological
structure in sequential settings.

Every problem has a modeled \textbf{intrinsic Betti number} : the number
of independent parallel paths that the problem's structure supports in
this abstraction. A blood test, an MRI, and a genetic screen are
\emph{diagnostically} independent --- each tests a different modality
(biochemistry, anatomy, genomics) and produces non-redundant
information, giving . The reflects that independence is a function of
the diagnostic question: for some conditions, a genetic result might
obviate the MRI (reducing ), while for others, all three are genuinely
independent. Eight compression codecs applied to the same chunk are
independent -- . The paths in a Feynman path integral are independent --
. In this framework, is estimated from the dependency structure rather
than treated as a direct design knob.

Every implementation has an \textbf{actual Betti number} : the number of
independent parallel paths in the system as built. A sequential referral
chain has . A fork with 8 codecs has . The gap between and is the
\textbf{topological deficit}:

When , the system's topology matches the problem's topology. It is
operating at its natural parallelism. When , the system is forcing a
high- problem through a low- pipe. The deficit is wasted parallelism --
performance left on the table.

I define the unit of topological deficit as the \textbf{Bule} (symbol:
\textbf{B}). One Bule equals one unit of -- one independent parallel
path that the problem supports but the implementation does not exploit.

A system at 0 B is topology-matched in this metric. A system at 3 B is
leaving three independent parallel paths unexploited. The Bule is
dimensionless, integer-valued in this representation, and computable
once a dependency graph and modeling protocol are specified.

This also yields an interpretive notion of \textbf{computational
aesthetics}: elegance corresponds to fit between implemented topology
and problem topology. In this framing, Bules are a structural diagnostic
of that fit.

The term is model-estimated rather than uniquely observer-independent in
open systems. Reported deficits should therefore be interpreted with
explicit modeling assumptions and uncertainty intervals where available.

\textbf{Topological deficit as a candidate diagnostic for real-world
waste.}

\textbf{System} \textbar{} \textbf{β1∗\hspace{0pt}} \textbar{}
\textbf{β1\hspace{0pt}} \textbar{} \textbf{Deficit} \textbar{}
\textbf{Observable Waste} \textbar{}

Healthcare diagnosis \textbar{} \textbar{} 0 (referral chain) \textbar{}
3 B \textbar{} 5-year average diagnosis time in the 2024 EURORDIS Rare
Barometer survey

Financial settlement \textbar{} 2 \textbar{} 0 (T+2 sequential)
\textbar{} 2 B \textbar{} order-of-\$4.4T lockup from a 2-day heuristic
applied to the DTCC/NSCC daily baseline; larger scenarios are
companion-model outputs

HTTP/2 multiplexing \textbar{} \textbar{} 0 (TCP substrate) \textbar{} B
\textbar{} Head-of-line blocking on any packet loss \textbar{}

Photosynthetic antenna \textbar{} (pigments) \textbar{} (quantum
coherence) \textbar{} 0 B \textbar{} high step-level energy-transfer
efficiency \textbar{}

Path integral \textbar{} \textbar{} \textbar{} 0 B \textbar{} Exact
quantum-mechanical predictions \textbar{}

DNA replication \textbar{} 1 (lagging strand) \textbar{} 1 (Okazaki
fragments) \textbar{} 0 B \textbar{} Replication matches leading strand
speed \textbar{}

Saltatory conduction \textbar{} nodes \textbar{} nodes \textbar{} 0 B
\textbar{} 100x speedup vs.~continuous conduction \textbar{}

In this paper's analyzed set, the pattern is: \textbf{cases coincide
with high-fit outcomes, and} \textbf{cases coincide with measurable
waste.} This is correlational evidence in the analyzed set, not a
standalone causal identification claim. The deficit is not abstract --
it maps to years of diagnostic delay, trillions of locked capital, and
protocol-level blocking.

This yields a practical diagnostic tool:

\begin{enumerate}
\def\labelenumi{\arabic{enumi}.}
\tightlist
\item
  \textbf{Measure} : analyze the problem's dependency structure to find
  its intrinsic parallelism. Independent inputs are independent paths.
  Sequential dependencies are constraints that reduce .\\
\item
  \textbf{Measure} : count the actual parallel paths in the
  implementation. A sequential pipeline has . A fork with paths has .\\
\item
  \textbf{Compute} : the gap is the optimization opportunity. If , the
  system is topologically constrained; micro-optimization within that
  fixed topology cannot recover the parallel paths the topology itself
  suppresses.
\end{enumerate}

\textbf{A scoped converse.} When a system exhibits fork/race/fold with
-- for example, photosynthesis, DNA replication and myelinated
conduction in this analyzed set -- that is evidence consistent with
near-optimal topological fit under this paper's constraints. It is not a
universal proof of unique optimality.

This is why the biological examples in §1 are not decoration. They are
supporting evidence for the correspondence hypothesis used here. When
\emph{Physarum} constructs transport networks with tradeoffs similar to
the Tokyo rail system, we observe a small measured deficit without
centralized design. When photosynthetic antenna complexes exhibit high
step-level transfer efficiency, we observe a high-fit topology for that
step of the process. These are selected examples, not universal proofs.

The optimality diagnostic also clarifies \textbf{one route to quantum
speedup}. Classical implementations with can carry a topological deficit
that quantum systems partially close by exploring paths concurrently.
For some problem families (for example, unstructured search), this
manifests as the familiar Grover-style gap

. But the converse does not hold in general: high structural readiness
does \emph{not} automatically imply Grover/Shor-style asymptotics

. In exact full-aggregation workloads (checksums, exact sums, full
histograms), the black-box cost still scales as because every item must
be read. In this paper's framing, is a structural feature worth
investigating, not a sufficient or theorem-like certificate of
asymptotic improvement.

In this framing, algorithmic aesthetics is an interpretive overlay on
measured topology mismatch. In the analyzed case studies, higher
deficits co-occur with years of diagnostic delay, large settlement
lockup, and protocol-level blocking. This is correlational evidence, not
standalone causal attribution.

\subsubsection{\texorpdfstring{\textbf{6.14 Map/Reduce as a
Topology-Readiness Screening Heuristic (Not a
Theorem)}}{6.14 Map/Reduce as a Topology-Readiness Screening Heuristic (Not a Theorem)}}\label{mapreduce-as-a-topology-readiness-screening-heuristic-not-a-theorem}

Map/reduce should be interpreted topologically (in the sense of the
MapReduce computation model

):

\begin{itemize}
\tightlist
\item
  \textbf{Map} is fork over independent partitions.\\
\item
  \textbf{Reduce} is fold under an associative/deterministic merger.\\
\item
  \textbf{Shuffle} is the routing layer between the two.
\end{itemize}

In this sense, map/reduce is a constrained fork/fold system with no
explicit race or vent semantics. That constraint is partly why
map/reduce usage is a useful signal: it usually means the workload
already has an exposed parallel frontier and a valid fold boundary. The
lack of race and vent semantics also indicates potentially unexploited
parallelism.

This leads to a practical claim:

\textbf{Heuristic claim.} Sustained map/reduce usage is evidence of
\textbf{topology readiness} for Wallington pipelines (fork/race/fold +
vent), and can motivate quantum-style path-exploration experiments in a
narrow structural sense (the problem admits concurrent path exploration
and deterministic projection).

This is \textbf{not} a claim of automatic quantum advantage. It does
\textbf{not} imply Grover/Shor-style asymptotics

. It only claims structural compatibility.

\textbf{Scoped heuristic.} Within the black-box workload simulations
used here, topology readiness is a useful screen for workloads worth
testing for quantum-style path exploration: without an exposed parallel
frontier, a deterministic fold boundary, and nonzero topological
opportunity, the companion model produces little or no migration gain

. Passing that screen is still not sufficient; asymptotic gain remains
family-dependent

.

\begin{description}
\tightlist
\item[Executable companion coverage makes this boundary explicit]
\hfill
\begin{enumerate}
\def\labelenumi{(\roman{enumi})}
\tightlist
\item
  when , modeled migration gain collapses to near-zero even with high
  map/reduce quality, (ii) a high- workload can still have no asymptotic
  quantum speedup (full aggregation: classical , quantum ), and (iii)
  another high-readiness family can exhibit Grover-style scaling
  (unstructured search: classical , quantum ). The heuristic therefore
  screens for topology compatibility, not algorithmic complexity class.
\end{enumerate}
\end{description}

I separate readiness from opportunity:

where all factors are normalized to :

\begin{itemize}
\tightlist
\item
  : fraction of map work that is truly independent (measured as the
  ratio of map tasks with zero cross-partition data access).\\
\item
  : reducer associativity/determinism score (1.0 if the reducer is
  associative and commutative; penalized for order dependence or
  non-determinism).\\
\item
  : partition skew (Gini coefficient of key distribution; 0 = uniform, 1
  = all keys in one partition).\\
\item
  : zero-copy ratio across map/shuffle/fold boundaries (fraction of data
  transferred without serialization).\\
\item
  : topological opportunity from the Bule deficit.
\end{itemize}

\textbf{Caveat.} These five factors are \emph{not} provably independent
--- and are likely correlated (high skew implies uneven independence),
and may constrain (non-associative reducers often require intermediate
serialization). The formula is a screening heuristic, not a calibrated
model. No threshold values are established for ``high'' vs.~``low'' ;
use here is ordinal (rank systems by , prioritize the highest for
experimental follow-up). It is not a standalone go/no-go rule.

Interpretation:

\begin{itemize}
\tightlist
\item
  High , low : architecture is ready, but little headroom (already near
  ).\\
\item
  Low , high : headroom exists, but map/reduce quality is too poor to
  realize it safely.\\
\item
  High : prioritize for preregistered pilot evaluation before migration
  to full Wallington primitives (add race + vent + Reynolds-driven
  multiplexing).
\end{itemize}

So map/reduce can be interpreted as a \textbf{screening diagnostic}: it
flags workloads likely to benefit from promotion into fork/race/fold,
and in a subset of cases may coincide with hypotheses worth testing for
quantum-style gains. The value is triage: it prioritizes which workloads
to test first. The formula has guided three internal production
migrations in the author's own systems (inference routing, session
preloading, and deploy artifact streaming --- all described in this
paper), but has not been independently validated beyond these cases and
should be treated as hypothesis-generating. An open-source
@affectively/aeon-pipelines implementation is available

.

\paragraph{\texorpdfstring{\textbf{Executable Diagnostic
Tool}}{Executable Diagnostic Tool}}\label{executable-diagnostic-tool}

The topological deficit is not just a theoretical quantity. The
@affectively/aeon package

includes a TopologyAnalyzer that computes Betti numbers and Bules from a
computation graph, and a TopologySampler that instruments a running
system to measure deficit over time:

import \{ TopologyAnalyzer, TopologySampler \} from `@affectively/aeon';

// Static analysis: is this system wasting parallelism?\\
const graph = TopologyAnalyzer.fromForkRaceFold(\{\\
forkWidth: 1, // implementation: sequential\\
intrinsicBeta1: 7, // problem: 8 independent codecs\\
\});\\
const report = TopologyAnalyzer.analyze(graph);\\
// → deficit: 7 Bules -- ``Sequential bottleneck: 7 Bules of waste''

// Fix: match the topology\\
const fixed = TopologyAnalyzer.fromForkRaceFold(\{\\
forkWidth: 8, // implementation: fork 8 codecs\\
intrinsicBeta1: 7, // problem: 8 independent codecs\\
\});\\
const fixedReport = TopologyAnalyzer.analyze(fixed);\\
// → deficit: 0 Bules -- ``Optimal: 0 Bules''

// Runtime sampling: how does the deficit evolve?\\
const sampler = new TopologySampler(\{ intrinsicBeta1: 7 \});\\
sampler.fork(`chunk-1', {[}`raw', `rle', `delta', `lz77',\\
`brotli', `gzip', `huffman', `dict'{]});\\
// → currentDeficit(): 0 Bules (all 8 codecs racing)\\
sampler.race(`chunk-1', `brotli');\\
sampler.vent(`chunk-1', `raw');\\
// \ldots{} vent remaining losers \ldots{}\\
sampler.fold(`chunk-1');\\
const samplerReport = sampler.report();\\
// → peakBeta1: 7, efficiency: 0.125 (1 race / 8 events)

The TopologyAnalyzer computes , , and detects fork/join pairs from any
directed graph. The TopologySampler records fork/race/vent/fold events
at runtime and produces time-series utilization data. Both are validated
by targeted tests covering sequential pipelines, fork/join graphs, void
detection, deficit measurement, concurrent forks, vent ratios, and the
real-world topologies from this section

.

\textbf{To see fork/race/fold is to see a system that has found its
shape. To not see it is to see a system that hasn't. The Bule count
tells you how far off you are.}

\subsection{\texorpdfstring{\textbf{7. Instantiation C: Distributed
Staged Computation (Stack Layer
3)}}{7. Instantiation C: Distributed Staged Computation (Stack Layer 3)}}\label{instantiation-c-distributed-staged-computation-stack-layer-3}

I implement fork/race/fold in a distributed computation engine with
processing stages partitioned across networked nodes -- a domain of
particular interest to the researcher. In the geometric framing, this
layer is where failure stops being a postmortem label and becomes a
measurable drift field.

In Gnosis (§11), the Wallington Rotation for a 4-stage pipeline is:

(tokens: Source \{ data: `workload' \})\\
(stage\_1: Node \{ id: `1' \}) (stage\_2: Node \{ id: `2' \})\\
(stage\_3: Node \{ id: `3' \}) (stage\_4: Node \{ id: `4' \})\\
(tokens)-{[}:FORK{]}-\textgreater(stage\_1 \textbar{} stage\_2
\textbar{} stage\_3 \textbar{} stage\_4)\\
(stage\_1 \textbar{} stage\_2 \textbar{} stage\_3 \textbar{}
stage\_4)-{[}:FOLD \{ strategy: `merge-all' \}{]}-\textgreater(result)

The topology is the program. The scheduling is the shape.

\subsubsection{\texorpdfstring{\textbf{7.1 Chunked Pipelined Prefill
(Wallington
Rotation)}}{7.1 Chunked Pipelined Prefill (Wallington Rotation)}}\label{chunked-pipelined-prefill-wallington-rotation}

In the baseline, a workload of items is processed sequentially through
stage nodes: round-trips. The key insight: each node's forward pass for
item depends only on that node's accumulated state from -- a stage-local
constraint (C1). This enables pipelining. Chunking groups items per
forward pass via causal masking.

The table below reports modeled step-count speedups only (not wall-clock
throughput), under A1-A2 above.

\textbf{Scenario} \textbar{} \textbf{Serial (P×N)} \textbar{}
\textbf{Chunked Pipeline} \textbar{} \textbf{Modeled Step-Count Speedup}
\textbar{}

14 tokens, 2 nodes \textbar{} 28 steps \textbar{} 9 steps \textbar{}
3.1x \textbar{}

100 tokens, 4 nodes \textbar{} 400 steps \textbar{} 7 steps \textbar{}
57x \textbar{}

500 tokens, 8 nodes \textbar{} 4,000 steps \textbar{} 15 steps
\textbar{} 267x \textbar{}

100 tokens, 10 nodes \textbar{} 1,000 steps \textbar{} 19 steps
\textbar{} 53x \textbar{}

\textbf{Measurement methodology.} Speedup figures are \emph{step-count
ratios} computed from the formula --- they measure scheduling depth
(number of sequential time steps), not wall-clock latency. Each ``step''
represents one chunk-stage processing event; per-step latency varies by
workload and hardware. The figures assume uniform stage latency and zero
inter-node communication cost (the benchmark harnesses mock network
communication, as noted in §13). Chunk size with chosen to maximize
throughput per the formula. These are \emph{theoretical best-case}
speedups for the scheduling topology; real-world figures would be
reduced by network RTT, uneven stage latencies, and queuing at node
boundaries. The 267x figure for 500 tokens / 8 nodes uses (one chunk),
giving steps.

In geometric terms, this is a stability law as much as a speed law.
Every reduction in exposed sequential depth shrinks the interval over
which outward drift can accumulate before the computation returns to its
sink. The Wallington Rotation reduces not only latency but
\textbf{failure exposure}.

\textbf{Wall-clock matrix evidence (fixture-scoped).} A live distributed
wall-clock matrix is provided via
companion-tests/scripts/gate1-wallclock-matrix.ts, with artifacts in
companion-tests/artifacts/gate1-wallclock-matrix.\{json,md\}. The
harness runs real loopback HTTP stage servers across predeclared
RTT/jitter/loss/workload cells, reporting p50/p95 completion latency
plus 95\% bootstrap confidence intervals and explicit pass/fail
criteria. In this matrix, all predeclared primary cells reject
no-improvement (speedup CI lower bound \textgreater{} 1.0 and
improvement CI lower bound \textgreater{} 0 ms). Non-loopback runs also
satisfy the same criteria in
companion-tests/artifacts/gate1-wallclock-external-single-host.\{json,md\}
and
companion-tests/artifacts/gate1-wallclock-external-multihost.\{json,md\}
(six distinct external hosts, one stage endpoint per host). This
supports a scoped wall-clock claim for this harness family and does not
by itself imply universal production-network speedups.

\subsubsection{\texorpdfstring{\textbf{7.2 Turbulent
Multiplexing}}{7.2 Turbulent Multiplexing}}\label{turbulent-multiplexing}

In molecular biology, a polysome (also called a polyribosome) is a
cluster of multiple ribosomes that are simultaneously translating a
single mRNA strand into proteins.

Think of it as a molecular assembly line: instead of one worker
(ribosome) reading an instruction manual (mRNA) and finishing the
product before the next one starts, multiple workers jump on the manual
as soon as the first one moves out of the way. This allows the cell to
mass-produce proteins with high throughput and efficiency.

When , 43 percent of node-slots are idle during ramp-up/ramp-down.
Turbulent multiplexing fills idle slots with items from concurrent
requests, maintaining per-request vent isolation (C2). This is analogous
to polysome behavior: fill the mRNA pipeline with multiple ribosomes,
degrade the mRNA when drops below threshold, and reallocate to active
pipelines. In the present framing, multiplexing is \textbf{fail-flow}:
when one request loses inward pressure, the manifold borrows stability
from other requests rather than letting the whole surface sag outward.

\subsubsection{\texorpdfstring{\textbf{7.3 Worthington Whip
(Superposition
Sharding)}}{7.3 Worthington Whip (Superposition Sharding)}}\label{worthington-whip-superposition-sharding}

A single workload is sharded across parallel pipelines. Each shard
processes items, then cross-shard correction reconciles at fold.
Per-shard compute savings: . The correction phase is the formal payment
demanded by our \texttt{THM-FAIL-MINCOST} logic: the manifold widened to
gain throughput, so it must later pay reconciliation tension to become
single-valued again.

\subsubsection{\texorpdfstring{\textbf{7.4 Speculative
Tree}}{7.4 Speculative Tree}}\label{speculative-tree}

A lightweight predictor generates candidate continuations (fork). All
branches enter the pipeline as multiplexed sub-requests (race). A
verifier checks all in a single batched pass. Invalid branches are
pruned via venting. Expected items accepted per pass with acceptance
rate : .

\subsection{\texorpdfstring{\textbf{8. Instantiation D: Aeon Flow
Protocol (Stack Layer
4)}}{8. Instantiation D: Aeon Flow Protocol (Stack Layer 4)}}\label{instantiation-d-aeon-flow-protocol-stack-layer-4}

\subsubsection{\texorpdfstring{\textbf{8.1 Design
Principle}}{8.1 Design Principle}}\label{design-principle}

The patterns -- fork, race, fold, vent -- recur with the same primitive
structure in edge composition, service worker preloading, fragment
assembly, deploy artifact streaming, CRDT synchronization and other
independent domains validated in §13. Rather than reimplementing per
domain, I extract the primitive into a binary wire protocol on UDP
dubbed Aeon Flow.

In Gnosis (§11), a multiplexed site load over Aeon Flow is:

(html: Asset \{ type: `text/html' \}) (css: Asset \{ type: `text/css'
\})\\
(js: Asset \{ type: `application/javascript' \}) (font: Asset \{ type:
`font/woff2' \})\\
(site)-{[}:FORK{]}-\textgreater(html \textbar{} css \textbar{} js
\textbar{} font)\\
(html \textbar{} css \textbar{} js \textbar{} font)-{[}:FOLD \{
strategy: `merge-all' \}{]}-\textgreater(cached\_site)

Four assets, one connection, one fold. The GGL program compiles directly
to the FlowFrame binary format below.

\subsubsection{\texorpdfstring{\textbf{8.2 Wire
Format}}{8.2 Wire Format}}\label{wire-format}

Offset Size Field\\
{[}0..1{]} u16 stream\_id (multiplexed stream identifier)\\
{[}2..5{]} u32 sequence (position within stream)\\
{[}6{]} u8 flags (FORK=0x01 \textbar{} RACE=0x02 \textbar{} FOLD=0x04
\textbar{} VENT=0x08 \textbar{} FIN=0x10)\\
{[}7..9{]} u24 length (payload bytes, max 16 MB)\\
{[}10..{]} {[}u8{]} payload (zerocopy Uint8Array view)

\textbf{10 bytes.} Every frame carries its own identity. Every frame is
self-describing. No ordered delivery is required. The stream\_id +
sequence pair is the coordinate in the covering space (§3.3). Flags
compose: RACE \textbar{} FIN means ``racing AND final frame.'' The frame
reassembler (§3.3) is the covering map back to sequential order.
Payloads are zerocopy: the codec writes 10 bytes in front of the
existing ArrayBuffer view. The header is therefore not mere metadata. It
is the \textbf{coordinate chart} of the manifold on which failure and
recovery are tracked.

\subsubsection{\texorpdfstring{\textbf{8.2.1 The Self-Describing Frame
as Pervasive
Abstraction}}{8.2.1 The Self-Describing Frame as Pervasive Abstraction}}\label{the-self-describing-frame-as-pervasive-abstraction}

The self-describing frame is not specific to the wire protocol. It is
the unifying data structure across both the transport layer and the
computation engine.

On the wire, it is the \textbf{FlowFrame} -- 10 bytes of header carrying
stream\_id, sequence, flags and length. On the computation side, it is
the \textbf{WorkFrame} -- the same (stream\_id, sequence) identity
enriched with a typed payload T and metadata:

{\def\LTcaptype{none} % do not increment counter
\begin{longtable}[]{@{}lr@{}}
\toprule\noalign{}
WorkFrame\textless T\textgreater{} & FlowFrame \\
\midrule\noalign{}
\endhead
\bottomrule\noalign{}
\endlastfoot
streamId: StreamId & streamId: u16 \\
sequence: number & sequence: u32 \\
payload: T & flags: u8 \\
metadata: Record\textless string,unknown\textgreater{} & length: u24 \\
emittedAt: number & payload: {[}u8{]} \\
\end{longtable}
}

The two are isomorphic, meaning they are the same shape with different
labels. The wire format bridge encodes
WorkFrame\textless T\textgreater{} to FlowFrame (serializing T as
payload bytes) and decodes FlowFrame back to
WorkFrame\textless T\textgreater. A computation that forks 10 streams
in-process produces 10 WorkFrames. Those same frames, encoded as
FlowFrames, can cross a network boundary and be reassembled on the other
side by the same FrameReassembler algorithm. The computation topology is
independent of the transport topology.

This is the same pattern as Okazaki fragments in DNA replication, chosen
to underscore the natural cohesion of this protocol design: each
fragment carries its genomic coordinate (its stream\_id + sequence),
enabling out-of-order synthesis and reassembly by DNA ligase. The
fragment is self-describing whether it is being synthesized on the
lagging strand (in-process) or transported via a virus to another cell
(on the wire). Identity is intrinsic, not assigned by context.

\subsubsection{\texorpdfstring{\textbf{8.3 Why UDP
Only}}{8.3 Why UDP Only}}\label{why-udp-only}

TCP had a long and successful run. For workloads with high
concurrent-path structure (), some TCP guarantees become tradeoffs:

\textbf{TCP Guarantee} \textbar{} \textbf{Why It Hurts} \textbar{}

Ordered delivery \textbar{} One lost packet on stream A blocks
\emph{all} streams behind it \textbar{}

Connection handshake \textbar{} 1.5 RTT before first data byte
\textbar{}

Single-stream congestion \textbar{} TCP backs off the entire connection
on loss \textbar{}

Connection-level retransmit \textbar{} Stream A's retransmit delays
stream B \textbar{}

HTTP/2 tried to multiplex streams over TCP. The application topology ()
contradicts the transport topology (). Head-of-line blocking is the
topological symptom (§3.4): two independent geodesics are forced to
share one ordered dimension. HTTP/3 (QUIC) partially resolves this with
per-stream loss recovery on UDP, but maintains ordered delivery within
each stream and retains a more complex framing surface than Aeon Flow in
this benchmark scope.

Aeon Flow -- a UDP-native alternative in this paper's benchmark scope --
does not patch TCP's problems at the application layer; it changes the
transport assumptions directly.

It starts from the topology and asks which wire format better fits
workloads: self-describing frames with no ordered delivery, AIMD
congestion control per-stream (not per-connection), MTU-aware
fragmentation (4-byte fragment header, 255 fragments × 1,468 bytes), and
ACK bitmaps (14 bytes covering 64 sequences). The protocol is about 800
lines of TypeScript. In the shootoff benchmarks used here, it
outperforms HTTP/3 on measured framing metrics and selected latency
measurements. These are benchmark-scoped results, not a universal
internet-wide claim; the topological-fit interpretation is a mechanism
hypothesis supported by these measurements

.

\subsubsection{\texorpdfstring{\textbf{8.4 Protocol
Comparison}}{8.4 Protocol Comparison}}\label{protocol-comparison}

\textbf{Metric} \textbar{} \textbf{HTTP/1.1} \textbar{} \textbf{HTTP/2}
\textbar{} \textbf{HTTP/3 (QUIC)} \textbar{} \textbf{Aeon Flow}
\textbar{}

Per-resource overhead \textbar{} \textasciitilde200 bytes \textbar{}
\textasciitilde30 bytes (HPACK) \textbar{} \textasciitilde20 bytes
\textbar{} \textbf{10 bytes} \textbar{}

Header fraction (12 resources) \textbar{} 4.2 percent \textbar{} 0.6
percent \textbar{} 0.4 percent \textbar{} \textbf{0.2 percent}
\textbar{}

Header fraction (95 resources) \textbar{} 31.0 percent \textbar{} 5.8
percent \textbar{} 4.1 percent \textbar{} \textbf{1.5 percent}
\textbar{}

Connections for full site \textbar{} 6+ \textbar{} 1 \textbar{} 1
\textbar{} \textbf{1} \textbar{}

Head-of-line blocking \textbar{} Yes (conn) \textbar{} Yes (TCP)
\textbar{} No (per-stream) \textbar{} \textbf{No} \textbar{}

Native fork/race/fold \textbar{} No \textbar{} No \textbar{} No
\textbar{} \textbf{Yes} \textbar{}

Vent propagation \textbar{} N/A \textbar{} RST\_STREAM \textbar{}
STOP\_SENDING \textbar{} \textbf{Recursive tree} \textbar{}

Transport \textbar{} TCP \textbar{} TCP \textbar{} UDP (QUIC) \textbar{}
\textbf{UDP (raw)} \textbar{}

Ordered delivery \textbar{} Required \textbar{} Required \textbar{}
Per-stream \textbar{} \textbf{None} \textbar{}

Topological contradiction \textbar{} N/A \textbar{} mismatch \textbar{}
Partial \textbar{} \textbf{None} \textbar{}

\subsubsection{\texorpdfstring{\textbf{8.5 Shootoff: Head-to-Head
Protocol
Benchmarks}}{8.5 Shootoff: Head-to-Head Protocol Benchmarks}}\label{shootoff-head-to-head-protocol-benchmarks}

I benchmark Aeon Flow against HTTP/1.1, HTTP/2 and HTTP/3 with realistic
compression (gzip, brotli) across two site profiles. All protocols use
identical payloads; only framing and transport differ. These are
deterministic fixture benchmarks for the specified payloads and
settings, not population-level estimates.

\textbf{Big Content Site} (12 resources, \textasciitilde2.5 MB -- large
JS bundles, hero images, web fonts):

\textbf{Protocol} \textbar{} \textbf{Wire Size} \textbar{}
\textbf{Framing Overhead} \textbar{} \textbf{Overhead \%} \textbar{}
\textbf{RTTs} \textbar{}

HTTP/1.1 \textbar{} 913 KB \textbar{} 8.2 KB \textbar{} 0.89 percent
\textbar{} 3 \textbar{}

HTTP/2 \textbar{} 907 KB \textbar{} 1.6 KB \textbar{} 0.18 percent
\textbar{} 2 \textbar{}

HTTP/3 (QUIC) \textbar{} 906 KB \textbar{} 906 B \textbar{} 0.10 percent
\textbar{} 1 \textbar{}

\textbf{Aeon Flow} \textbar{} \textbf{905 KB} \textbar{} \textbf{276 B}
\textbar{} \textbf{0.03 percent} \textbar{} \textbf{1} \textbar{}

For large payloads in this benchmark set, protocol wire sizes are close
-- but Aeon Flow's framing is \textbf{3.3x smaller than HTTP/3} (276 B
vs 906 B).

\textbf{Microfrontend Site} (95 resources, \textasciitilde1.8 MB -- 45
JS modules, 16 CSS modules, 20 SVG icons):

\textbf{Protocol} \textbar{} \textbf{Wire Size} \textbar{}
\textbf{Framing Overhead} \textbar{} \textbf{Overhead \%} \textbar{}
\textbf{RTTs} \textbar{}

HTTP/1.1 \textbar{} 187 KB \textbar{} 58.1 KB \textbar{} \textbf{31.0
percent} \textbar{} 16 \textbar{}

HTTP/2 \textbar{} 137 KB \textbar{} 8.0 KB \textbar{} 5.8 percent
\textbar{} 2 \textbar{}

HTTP/3 (QUIC) \textbar{} 135 KB \textbar{} 5.9 KB \textbar{} 4.4 percent
\textbar{} 1 \textbar{}

\textbf{Aeon Flow} \textbar{} \textbf{131 KB} \textbar{} \textbf{1.9 KB}
\textbar{} \textbf{1.5 percent} \textbar{} \textbf{1} \textbar{}

In this benchmark fixture, topology appears to matter. HTTP/1.1 wastes
\textbf{31 percent of total bandwidth on headers} -- nearly a third of
the wire is framing, not data. HTTP/2 reduces this to 5.8 percent.
HTTP/3 to 4.4 percent. Aeon Flow: \textbf{1.5 percent}. That is a
\textbf{21x reduction} in framing overhead versus HTTP/1.1 and
\textbf{3x versus HTTP/3} for this case.

At an illustrative 100ms RTT (ignoring loss/retransmit dynamics),
HTTP/1.1's 16 round trips imply \textasciitilde1.6 seconds of round-trip
latency budget. Aeon Flow: 1 round trip (\textasciitilde0.1 seconds on
the same RTT assumption). This gap is consistent with a topological
interpretation: HTTP/1.1 has (one request per connection, six
connections). Aeon Flow has (95 streams, one connection). The framing
overhead here is consistent with forcing a high- problem through a low-
pipe.

Modern frontend workloads often ship many small assets after
tree-shaking and code splitting, which amplifies request/metadata
overhead. In this benchmark scope, Aeon Flow multiplexes these assets
through one transport session and reduces framing cost. Effects on CLS,
INP and hydration strategy remain application-dependent and are not
guaranteed by transport alone.

\subsection{\texorpdfstring{\textbf{9. Instantiation E: Topological
Compression (Stack Layer 5 ---
Capstone)}}{9. Instantiation E: Topological Compression (Stack Layer 5 --- Capstone)}}\label{instantiation-e-topological-compression-stack-layer-5-capstone}

\subsubsection{\texorpdfstring{\textbf{9.1 The Claim and Its
Limits}}{9.1 The Claim and Its Limits}}\label{the-claim-and-its-limits}

The same fork/race/fold primitive applies to compression.
\textbf{Topological compression} forks all available codecs per chunk,
races them and folds to the winner. Each chunk independently selects its
best codec. The output is a sequence of self-describing frames (9-byte
header: codec ID, original size, compressed size). .

The geometric claim is not that compression has one universal champion.
It is that the two-level race maps a \textbf{Pareto frontier} between
ratio, portability, random access, and collapse cost. The interesting
question is therefore: what is the \textbf{minimal cost floor} at which
the system can return to one encoding while keeping the failure ledger
bounded?

In Gnosis (§11), the topological compressor is:

(raw: Codec \{ type: `raw' \}) (rle: Codec \{ type: `rle' \})\\
(brotli: Codec \{ type: `brotli' \}) (gzip: Codec \{ type: `gzip' \})\\
(chunk)-{[}:FORK{]}-\textgreater(raw \textbar{} rle \textbar{} brotli
\textbar{} gzip)\\
(raw \textbar{} rle \textbar{} brotli \textbar{}
gzip)-{[}:RACE{]}-\textgreater(smallest)

Four lines. The entire compression strategy. The topology is the
algorithm.

I implement this with eight codecs:

\textbf{ID} \textbar{} \textbf{Codec} \textbar{} \textbf{Type}
\textbar{} \textbf{Best on} \textbar{}

0 \textbar{} Raw (identity) \textbar{} Pure JS \textbar{} Incompressible
data \textbar{}

1 \textbar{} RLE \textbar{} Pure JS \textbar{} Repeated byte runs
\textbar{}

2 \textbar{} Delta \textbar{} Pure JS \textbar{} Sequential/incremental
data \textbar{}

3 \textbar{} LZ77 \textbar{} Pure JS \textbar{} Repeated patterns
\textbar{}

4 \textbar{} Brotli \textbar{} Platform (node:zlib) \textbar{} General
text \textbar{}

5 \textbar{} Gzip \textbar{} Platform (node:zlib) \textbar{} General
text (broad fallback) \textbar{}

6 \textbar{} Huffman \textbar{} Pure JS \textbar{} Skewed byte
distributions \textbar{}

7 \textbar{} Dictionary \textbar{} Pure JS \textbar{} Web content
(HTML/CSS/JS keywords) \textbar{}

Before any excitement takes hold, it is important to state a boundary:
fork/race/fold provides a container for adaptive codec selection, not a
guaranteed ratio improvement over the best standalone codec on
homogeneous payloads. Its value here is strategy selection,
composability and bounded framing overhead.

\subsubsection{\texorpdfstring{\textbf{9.2 What the Benchmarks Actually
Show}}{9.2 What the Benchmarks Actually Show}}\label{what-the-benchmarks-actually-show}

I benchmark across both sites on Aeon Flow transport. The results are
honest and fixture-specific:

\textbf{Big Content Site} (12 resources, \textasciitilde2.22 MB):

\textbf{Compression} \textbar{} \textbf{Wire Size} \textbar{}
\textbf{Ratio} \textbar{} \textbf{β1\hspace{0pt}} \textbar{}

Brotli (global, quality 4) \textbar{} 905 KB \textbar{} 39.8 percent
\textbar{} 0 \textbar{}

Topo-full (8 codecs per-chunk) \textbar{} 1005 KB \textbar{} 44.2
percent \textbar{} 7 \textbar{}

Topo-pure (6 pure-JS codecs per-chunk) \textbar{} 1.17 MB \textbar{}
52.5 percent \textbar{} 5 \textbar{}

\textbf{Microfrontend Site} (95 resources, \textasciitilde617 KB):

\textbf{Compression} \textbar{} \textbf{Wire Size} \textbar{}
\textbf{Ratio} \textbar{} \textbf{β1\hspace{0pt}} \textbar{}

Brotli (global, quality 4) \textbar{} 131 KB \textbar{} 20.9 percent
\textbar{} 0 \textbar{}

Topo-full (8 codecs per-chunk) \textbar{} 159 KB \textbar{} 25.4 percent
\textbar{} 7 \textbar{}

Topo-pure (6 pure-JS codecs per-chunk) \textbar{} 229 KB \textbar{} 36.8
percent \textbar{} 5 \textbar{}

\textbf{Standalone brotli wins on compression ratio.} On these
benchmarks -- homogeneous web content -- global brotli beats per-chunk
topological compression by 4--15 percentage points. This is not
surprising: brotli compresses the entire stream with a sliding window
that builds dictionary context across chunks. Per-chunk compression
resets the dictionary every 4096 bytes.

The two-level race (§9.3) confirms this. On these payloads, when given
the choice between global brotli and per-chunk topology, the
harness-selected winner was global brotli across the observed benchmark
runs, matching standalone brotli's ratio plus 5 bytes of strategy
header. For this homogeneous-content benchmark, the 9-byte per-chunk
header tax and the loss of cross-chunk dictionary context outweighed
per-chunk adaptive gains.

\subsubsection{\texorpdfstring{\textbf{9.3 Two-Level Stream
Race}}{9.3 Two-Level Stream Race}}\label{two-level-stream-race}

I extend the topology to race at two levels:

fork (stream level):\\
\textbar- Path 0: Per-chunk topological (8 codecs × each 4096-byte
chunk)\\
\textbar- Path 1: Global brotli (entire stream, cross-chunk
dictionary)\\
\textbar- Path 2: Global gzip (entire stream)\\
`- \ldots{}\\
race → smallest total output wins\\
fold → 5-byte strategy header + compressed data

This is the usefulness of fork/race/fold to compression: with brotli
included as a racing path, the stream-level strategy tracks brotli's
ratio within a bounded strategy-header overhead. On these benchmarks it
is not better than standalone brotli; the observed downside is the fixed
5-byte strategy header.

\subsubsection{\texorpdfstring{\textbf{9.4 What the Topology Actually
Provides}}{9.4 What the Topology Actually Provides}}\label{what-the-topology-actually-provides}

If topological compression does not beat brotli on ratio, what is the
point?

\textbf{1. Subsumption, not superiority.} The topology is the space in
which brotli competes. Brotli at is a degenerate case of topological
compression at . The two-level race includes brotli as a contestant. If
brotli is best, the topology selects it. If something better appears
tomorrow -- a learned codec, a neural compressor, a domain-specific
dictionary -- it enters the race without changing the architecture. The
TopologicalCompressor is unchanged; only the codec array grows. In
geometric terms, the system searches the frontier without presuming
where the sink lies.

\textbf{2. Platform independence.} Brotli requires node:zlib (Node, Bun,
Deno). In browsers and Cloudflare Workers, it is unavailable. Aeon's
\texttt{topo-pure} -- six codecs in pure JavaScript, zero dependencies
-- achieves 36.8 percent ratio on the microfrontend with no native code.
The topology degrades gracefully: full ratio when brotli is available,
reasonable ratio when it is not. For software engineers, there is
technical value in fewer dependencies. For people, it helps set the
table for a serverless ecosystem built on a local-first technology
stack.

\textbf{3. Per-chunk random access.} The per-chunk format enables
decompression of individual chunks without processing the entire stream.
For seeking into large payloads, resuming interrupted transfers, or
parallel decompression, monolithic global brotli requires external
indexing to provide comparable access.

\textbf{4. Adaptive codec selection on heterogeneous data.} On the
per-chunk level, different regions of the input genuinely select
different codecs. The shootoff shows 3 distinct codecs winning across
151 chunks on realistic web content (brotli for text chunks, dictionary
for web-pattern-heavy chunks, raw for incompressible binary). Within
this tested codec set and strategy surface, no single fixed codec
reproduces that per-chunk winner diversity.

\textbf{5. The real compression win is framing, not codecs.} The paper's
compression contribution is not beating brotli's ratio. In the
microfrontend benchmark, it is the 30× reduction in framing overhead
(§8.4): Aeon Flow uses 1.9 KB of framing for 95 resources where HTTP/1.1
uses 56.3 KB. On that fixture, framing overhead drops from 31.0 percent
to 1.5 percent of total wire bytes. This saving is orthogonal to which
codec compresses the content.

\subsubsection{\texorpdfstring{\textbf{9.5 Honest
Assessment}}{9.5 Honest Assessment}}\label{honest-assessment}

The per-chunk topological approach pays a real cost: 9 bytes per chunk
of header overhead and the loss of cross-chunk dictionary context. On
the homogeneous content used in this benchmark set, this cost exceeds
the benefit of adaptive codec selection. Global brotli, with its
full-stream dictionary, simply compresses text better than any per-chunk
approach can.

\textbf{Comparison to adaptive single-algorithm heuristics.} A simpler
alternative --- ``use brotli for text, raw for binary, based on
content-type heuristic'' --- would capture most of the per-chunk
topology's adaptive benefit at zero per-chunk overhead. On these
benchmarks, such a heuristic is expected to be close to global brotli's
ratio (because payloads are predominantly web text). The per-chunk
topology's advantage over simple heuristics emerges only on
\emph{heterogeneous} payloads (mixed binary/text, embedded images in
HTML, protocol buffers interleaved with JSON) where content-type
heuristics misclassify regions. The shootoff's 3-codec-winner
distribution across 151 chunks is an initial indication of this
behavior: even on mostly-homogeneous web content, 12 percent of chunks
selected a non-brotli winner (dictionary for web-pattern-heavy chunks,
raw for incompressible binary).

The two-level stream race eliminates this disadvantage by including
global brotli as a racing path. But it also reveals that per-chunk
topological compression, as implemented here, is not the winning
strategy for web content. It is a structurally sound framework that
provides platform independence, random access and future extensibility
-- at the cost of matching, not beating, the state of the art on ratio.

The progression four codecs () → six codecs () → eight codecs ()
demonstrates the covering-space property: each expansion improved
pure-JS compression without changing the base space. But adding brotli
and gzip to the race, while improving per-chunk results, did not
overcome the global-dictionary advantage on these benchmarked workloads.

\textbf{The topological framework subsumes individual codec strategies.
It does not necessarily surpass the best one on ratio.} On the evaluated
web-content workloads, topological compression with per-chunk racing did
not outperform global brotli ratio. Global brotli's full-stream
dictionary context retained a strong information advantage for these
inputs. The practical conclusion is that topology provides structural
guarantees -- strategy subsumption, platform independence, random
access, extensibility -- without guaranteeing ratio superiority on
homogeneous content.

Executable evidence is available in two independent suites: the
companion topological-compression obligations

and the production TopologicalCompressor tests in the open-source
@affectively/aeon package

. Together they verify per-chunk adaptive winner selection, 9-byte
self-describing chunk headers, codec vent behavior (discarding
expansions), two-level stream race strategy selection, invariants and
roundtrip correctness across edge cases and large payloads.

\subsubsection{\texorpdfstring{\textbf{9.6
Applications}}{9.6 Applications}}\label{applications}

\textbf{Aeon Foundation} \textbar{} \textbf{Fork} \textbar{}
\textbf{Race} \textbar{} \textbf{Fold} \textbar{}

\textbf{Site preloading} \textbar{} Stream all assets as parallel frames
\textbar{} First complete asset wins cache slot \textbar{} SW stores all
in Cache API \textbar{}

\textbf{Edge-Side Inference (ESI) composition} \textbar{} Fork stream
per directive \textbar{} Race cache vs.~compute \textbar{} Assemble into
final page \textbar{}

\textbf{Deploy artifacts} \textbar{} Fork per build artifact \textbar{}
Stream concurrently \textbar{} Receive complete deployment \textbar{}

\textbf{CRDT sync} \textbar{} Fork per-peer delta streams \textbar{}
Race peers to contribute \textbar{} Merge deltas into canonical state
\textbar{}

\textbf{Speculative nav} \textbar{} Fork predicted route preloads
\textbar{} Race prediction vs.~actual \textbar{} Display whichever
resolves first \textbar{}

\subsection{\texorpdfstring{\textbf{10. Instantiation A:
Self-Verification (Stack Layer 1 ---
Foundation)}}{10. Instantiation A: Self-Verification (Stack Layer 1 --- Foundation)}}\label{instantiation-a-self-verification-stack-layer-1-foundation}

A strong executable result for expressiveness in this scope: the model
checker can verify a model of its own exploration. In the geometric
presentation, that self-check becomes the prototype of FMFV: the checker
does not only count events, it witnesses its own return toward
stability.

\subsubsection{\texorpdfstring{\textbf{10.1 The Checker's BFS Is
Fork/Race/Fold}}{10.1 The Checker's BFS Is Fork/Race/Fold}}\label{the-checkers-bfs-is-forkracefold}

The ForkRaceFoldModelChecker in @affectively/aeon-logic

explores state spaces via breadth-first search. Each BFS layer is a time
step. Each state is a spatial position. The exploration graph maps
directly to the four primitives:

\textbf{BFS Operation} \textbar{} \textbf{Fork/Race/Fold Primitive}
\textbar{} \textbf{Topological Effect} \textbar{}

Expansion with \textgreater1 successor \textbar{} \textbf{Fork}
\textbar{} \textbar{}

Transition to already-visited state \textbar{} \textbf{Fold}
(interference) \textbar{} Creates independent cycle \textbar{}

Unfair cycle filtered by weak fairness \textbar{} \textbf{Vent}
\textbar{} Irreversible path removal \textbar{}

Frontier exhausted, exploration complete \textbar{} \textbf{Collapse}
\textbar{} \textbar{}

The checker computes and returns topological diagnostics
(CheckerTopologyStats) for every verification: forkCount, foldCount,
ventCount, beta1 (first Betti number of the exploration graph), and
depthLayers (path-integral time steps).

\subsubsection{\texorpdfstring{\textbf{10.2 Self-Verification as
TemporalModel}}{10.2 Self-Verification as TemporalModel}}\label{self-verification-as-temporalmodel}

The checker's own BFS exploration is modeled as a
TemporalModel\textless CheckerState\textgreater{} with 8 state variables
(explored, frontier, transitions, folded, forks, vents, depth, done) and
6 actions (ExpandLinear, ExpandFork, FoldTransition, VentCycle,
CompleteLayer, Finish). Another instance of the same checker verifies 7
invariants about this model:

\begin{enumerate}
\def\labelenumi{\arabic{enumi}.}
\tightlist
\item
  --- topology is well-formed\\
\item
  --- every back-edge creates exactly one independent cycle\\
\item
  --- you can only vent what has been folded\\
\item
  --- folds are a subset of transitions\\
\item
  --- at least the initial state\\
\item
  --- non-negative frontier\\
\item
  --- bounded exploration
\end{enumerate}

Liveness: (eventual termination) under weak fairness .

These boundary checks are packaged inside a \textbf{Foster-Lyapunov
Drift Schema}. Let the checker's Lyapunov potential be a simple
structural distance-to-sink measure, for example:

and let be the small set of termination-ready states with bounded
frontier and zero unresolved debt. The stability obligation is then:

In this presentation, the seven invariants are local boundary
conditions. The drift obligation is the global proof that the Void of
failure remains bounded.

\subsubsection{\texorpdfstring{\textbf{10.3 TLA+
Self-Verification}}{10.3 TLA+ Self-Verification}}\label{tla-self-verification}

The same model is rendered as a TLA+ specification via
renderSelfVerificationArtifactPair(), producing a .tla module (extending
Naturals, with weak fairness WF\_vars(Finish)) and a .cfg config. The
specification is validated through the runTlaSandbox() round-trip: parse
render parse identical. A dual verification test confirms both paths
agree: the TLA sandbox validates the spec structure, the checker
verifies the same model's invariants and liveness.

\subsubsection{\texorpdfstring{\textbf{10.4 Closure Under
Self-Application}}{10.4 Closure Under Self-Application}}\label{closure-under-self-application}

In the finite-model scope used here, self-verification provides a
constructive closure result. The topology stats the checker reports
about verifying itself (forkCount, foldCount, beta1) are themselves
fork/race/fold observables. The meta-topology -- the topology of the
checker checking itself -- has forks (multiple actions enabled per
state), folds (different action sequences reaching the same checker
state), and measurable .

This means fork/race/fold is closed under self-application: a system
built from these primitives can reason about systems built from these
primitives. The topological deficit of self-verification measures the
cost of self-knowledge, and the drift witness tells us whether that
self-knowledge still points back toward the sink.

Executable companion tests verify these claims

.

\subsection{\texorpdfstring{\textbf{11. Instantiation B: The Topological
Type-Checker (Stack Layer
2)}}{11. Instantiation B: The Topological Type-Checker (Stack Layer 2)}}\label{instantiation-b-the-topological-type-checker-stack-layer-2}

In classical programming languages, a type-checker verifies that a
string is not passed to an integer function. It protects memory layout.
In distributed systems, the most catastrophic failures are not memory
type errors; they are \textbf{topological expansions} -- queues that
grow without bound, forks that never fold, and retries that overwhelm
the network.

Gnosis

is a programming language that dispenses with imperative control flow to
make the computation graph explicit. More fundamentally, its compiler
(Betti) operates as a \textbf{topological auditor}. It does not merely
check syntax; it checks the geometry of the system. At this layer, the
thermodynamic identity acts as a compile-time accounting law: if the
declared topology has no stable geometric floor, the program is
ill-typed.

\subsubsection{\texorpdfstring{\textbf{11.1 From Syntax to Curvature:
The GGL State
Space}}{11.1 From Syntax to Curvature: The GGL State Space}}\label{from-syntax-to-curvature-the-ggl-state-space}

In Gnosis Graph Language (GGL), the source code defines the manifold. To
enforce stability, nodes and edges carry the physical forces acting on
that manifold: \textbf{pressure} (), \textbf{potential} (), and
\textbf{restorative drift} ().

Nodes define the potential landscape:

\begin{itemize}
\tightlist
\item
  :Source injects workload pressure.\\
\item
  :State defines the active manifold and its Lyapunov potential (for
  example, queue depth or count).\\
\item
  :Sink defines the stable equilibrium where .
\end{itemize}

Edges define the transition kernel and restorative forces. Rather than
simply routing data, FOLD and VENT edges can declare the coefficients
that make stability explicit:

(traffic: Source \{ pressure: `lambda' \})\\
(processing: State \{ potential: `beta1' \})\\
(complete: Sink \{ beta1\_target: 0 \})

// Fork injects potential energy\\
(traffic) -{[}:FORK \{ weight: 1.0 \}{]}-\textgreater{} (processing)

// Fold extracts useful work (service rate)\\
(processing) -{[}:FOLD \{\\
strategy: `winner-take-all',\\
service\_rate: `mu',\\
drift\_gamma: 1.0\\
\}{]}-\textgreater{} (complete)

// Vent provides the FMFV valve\\
(processing) -{[}:VENT \{\\
condition: `timeout \textgreater{} 100ms',\\
drift\_coefficient: `alpha(n)',\\
repair\_debt: 0\\
\}{]}-\textgreater{} (complete)

The language provides four core edge types --- FORK, RACE, FOLD, VENT
--- plus PROCESS and INTERFERE. The compiler reads the graph not just as
control flow, but as a \textbf{state space with curvature}.

\subsubsection{\texorpdfstring{\textbf{11.2 The Spectral Radius
Constraint}}{11.2 The Spectral Radius Constraint}}\label{the-spectral-radius-constraint}

When the Betti compiler parses this GGL code, it first extracts the
transition kernel from the declared edge weights, service rates, and
drift annotations.

Before a single byte of WASM is generated, Betti performs a
\textbf{spectral analysis}. It calculates the spectral radius of the
routing matrix. If , the program is geometrically unstable: topological
expansion exceeds the system's ability to collapse the state space back
toward a sink.

Betti rejects the program with ERR\_SPECTRAL\_EXPLOSION. The developer
is prevented from deploying a system that is mathematically predisposed
to tear the manifold open.

\subsubsection{\texorpdfstring{\textbf{11.3 Mechanized Drift and the
First
Law}}{11.3 Mechanized Drift and the First Law}}\label{mechanized-drift-and-the-first-law}

For state-dependent programs -- such as the adaptive alpha(n) vent
strategy above -- a static spectral check is insufficient. The
restorative force changes as failure debt grows.

To certify these topologies, Betti acts as a bridge to formal
verification. The compiler auto-generates a Lean 4 theorem representing
the graph's transition kernel and invokes a custom automated tactic,
derive\_gnosis\_drift.

That tactic mechanizes the Foster-Lyapunov drift schema. It isolates the
structural forces -- pressure versus restorative service and venting --
and calculates the expected change in potential:

The compiler then asks Lean to prove that curvature points strictly
inward:

where is a finite small set. If the developer writes a pipeline with
massive FORK pressure but no sufficient VENT edge, then remains zero.
The tactic sees . If pressure exceeds service capacity, Lean cannot
prove negative drift, and Betti halts compilation with
ERR\_DRIFT\_POSITIVE.

This is the topological type error. The program is rejected not because
it is syntactically malformed, but because it has no proven geometric
mechanism to dissipate excess potential when fold extraction cannot keep
up with fork injection.

\subsubsection{\texorpdfstring{\textbf{11.4 Runtime FMFV
Emission}}{11.4 Runtime FMFV Emission}}\label{runtime-fmfv-emission}

When the Lean tactic succeeds, the program is mathematically certified
to be geometrically stable. Betti then goes one step further.

Because the compiler used an explicit Foster-Lyapunov schema, it can
extract the certified bounds of the stable manifold -- the constants
(restorative pressure) and (the redline boundary of the small set).
Betti emits those verified constants alongside the generated FlowFrame
control tables and runtime metadata that accompany the WASM binary.

The compiled runtime possesses \textbf{Failure-Modulated Flow Validation
(FMFV)}. It does not guess when it is overwhelmed. It has a certified
operating envelope, allowing active nodes to trigger topological
contraction -- typically venting or load shedding -- exactly when the
state space crosses the geometric redline, preventing the manifold from
tearing.

The path TypeScript (Betty) -\textgreater{} GGL (Betti) -\textgreater{}
everything else establishes closure under construction for a compiler
that can certify its own stability before it emits the next program
image.

Gnosis supports a strong evidence-backed claim: it is a self-hosted,
self-checking topology language with automated formal-artifact
generation. The compiler topology is itself written in GG (betti.gg) and
included in formal lint checks, while execution paths enforce
bounded-state structural verification with explicit invariants and
eventual reachability conditions before or during topology use. The
verify workflow can generate TLC-ready TLA+ modules and configs with
safety and liveness obligations, and these paths are covered by
source-level tests and formal-check scripts

.

This is a claim of structural formal compatibility and mechanized
verification workflow, not a claim of automatic asymptotic quantum
advantage.

\subsubsection{\texorpdfstring{\textbf{11.5 The Five Domains as a
Stack}}{11.5 The Five Domains as a Stack}}\label{the-five-domains-as-a-stack}

The five instantiation domains are not independent --- they form a
stack, each enabled by the ones below:

\textbf{Stack Layer} \textbar{} \textbf{Domain} \textbar{} \textbf{§}
\textbar{} \textbf{Primitive} \textbar{} \textbf{Role} \textbar{}

1 (foundation) \textbar{} Self-verification \textbar{} §10 \textbar{}
Temporal model checking \textbar{} Verifies modeled invariants
\textbar{}

2 \textbar{} Topological type-checker \textbar{} §11 \textbar{} GGL +
Betty/Betti \textbar{} Compile-time stability audit \textbar{}

3 \textbar{} Distributed computation \textbar{} §7 \textbar{} Wallington
Rotation \textbar{} The scheduling algorithm \textbar{}

4 \textbar{} Edge transport \textbar{} §8 \textbar{} 10-byte FlowFrame
\textbar{} The wire format \textbar{}

5 (capstone) \textbar{} Compression \textbar{} §9 \textbar{} Per-chunk
codec racing \textbar{} Bytes on wire \textbar{}

The stack reads bottom-up: \emph{from building blocks to bytes on wire}.
Layer 1 (§10) verifies modeled primitive properties. Layer 2 (§11) gives
a language that type-checks topology and drift before deployment. Layer
3 (§7) schedules work through the topology, expressed in layer 2's
language. Layer 4 (§8) puts frames on the wire, carrying layer 3's
scheduled work. Layer 5 (§9) compresses the payload --- actual bytes,
actual ratios, actual wire --- using layers below it.

The Rust/WASM runtime executes the FlowFrames at the same byte-level
format defined in §8.2. The language is not a wrapper around the
protocol --- it is the protocol's native programming model.

The stack is the paper's clearest existence demonstration: one set of
four primitives (fork, race, fold, vent) yields a scheduling algorithm,
wire protocol, compression strategy, verification engine, and
programming language. Each layer is independently useful. Together they
form a computational ecosystem where topology, program structure, and
protocol design are aligned.

\subsection{\texorpdfstring{\textbf{12. The
Engine}}{12. The Engine}}\label{the-engine}

The algorithm is implemented as \textbf{Aeon Pipelines}

, a zero-dependency computation topology engine in TypeScript. It runs
on Cloudflare Workers, Deno, Node, Bun and browsers. The API surface is
two classes:

\begin{itemize}
\tightlist
\item
  \textbf{Pipeline}: the engine -- capacity, metrics, backpressure,
  turbulent multiplexing.\\
\item
  \textbf{Superposition\textless T\textgreater{}}: the builder --
  chainable fork/race/fold/vent/tunnel/interfere/entangle/measure/search
  operations.
\end{itemize}

// Kids juggling balls\\
const result = await Pipeline\\
.from({[}fetchFromA, fetchFromB, fetchFromC{]})\\
.race();

// People juggling the kids\\
const diagnosis = await Pipeline\\
.from({[}bloodTest, mriScan, geneticScreen{]})\\
.vent(result =\textgreater{} result.inconclusive)\\
.tunnel(result =\textgreater{} result.conclusive)\\
.fold(\{\\
type: `merge-all',\\
merge: mergeFindings,\\
\});

// Grover-style search over solution space\\
const drug = await new Pipeline(\{ capacity: 64 \})\\
.fork(candidates.map(c =\textgreater{} () =\textgreater{}
evaluate(c)))\\
.search(\{\\
width: 32,\\
oracle: compound =\textgreater{} compound.efficacy,\\
mutate: (compound, gen) =\textgreater{} perturb(compound, gen),\\
convergenceThreshold: 0.01,\\
\});

The search() operation is a classical heuristic inspired by Grover-style
amplification patterns. In some landscapes it reduces empirical
iteration counts versus naive sequential search, but this is not an
asymptotic complexity claim.

\subsubsection{\texorpdfstring{\textbf{12.1
Performance}}{12.1 Performance}}\label{performance}

The pipeline engine is designed for low orchestration overhead. In the
microbenchmarks below, orchestration cost is in the microsecond range,
and profiled workloads are typically dominated by user work functions.
These latency values are point estimates from the current
harness/environment and should be treated as order-of-magnitude
indicators rather than cross-machine constants.

A stronger statement is mechanized as a conditional formal obligation in
SchedulerBound.tla: under finite-topology execution with bounded frame
metadata and constant-time scheduler primitives, scheduler transition
cost is an additive bounded term independent of user-handler runtime.
This justifies ``handler-dominated runtime'' only within those explicit
assumptions, not as a universal claim

.

\textbf{Operation} \textbar{} \textbf{Latency} \textbar{} \textbf{Notes}
\textbar{}

fork(10) \textbar{} \textbf{1.82 µs} \textbar{} 10 parallel streams
created \textbar{}

fold(\{ type: `quorum', threshold: 3 \}) \textbar{} \textbf{4.51 µs}
\textbar{} Byzantine agreement across 5 streams \textbar{}

search(8×20) \textbar{} \textbf{8.3 µs} \textbar{} Grover-style search,
8-wide, 20 generations \textbar{}

interference(100) \textbar{} \textbf{16.3 µs} \textbar{} Pairwise
consensus across 100 streams \textbar{}

vent-tree(13) \textbar{} \textbf{18.9 µs} \textbar{} Recursive vent
across 13-node tree \textbar{}

flow-bridge-batch(100) \textbar{} \textbf{25.7 µs} \textbar{} 100 frames
encoded to wire format \textbar{}

reassemble-reverse(1000) \textbar{} \textbf{71.4 µs} \textbar{} 1,000
frames reassembled from reverse order \textbar{}

flow-bridge-roundtrip \textbar{} \textbf{0.76 µs} \textbar{} Single
frame encode → decode \textbar{}

Zero dependencies. \textasciitilde384 bytes per stream and
\textasciitilde3.5 KB per pipeline. Requires no servers.

\subsubsection{\texorpdfstring{\textbf{12.2 Domain
Validation}}{12.2 Domain Validation}}\label{domain-validation}

The same API -- unchanged -- was exercised in executable scenario
harnesses across multiple domain archetypes, including:

\begin{enumerate}
\def\labelenumi{\arabic{enumi}.}
\tightlist
\item
  \textbf{Multi-venue trading}: fork/race across exchanges, vent adverse
  prices\\
\item
  \textbf{Healthcare diagnostics}: fork parallel tests, tunnel on
  conclusive, merge-all findings\\
\item
  \textbf{Financial settlement}: fork clearing/netting/DVP, merge-all
  for T+0\\
\item
  \textbf{Construction scheduling}: fork trades per floor, merge-all
  hours\\
\item
  \textbf{Emergency dispatch}: fork/race responders, first arrival
  wins\\
\item
  \textbf{Academic review}: fork reviewers, quorum 2/3 agreement\\
\item
  \textbf{Drug discovery}: fork compounds, Grover search to
  convergence\\
\item
  \textbf{Manufacturing QC}: fork sensors, consensus (constructive
  interference)\\
\item
  \textbf{Journal publishing}: fork reviewers, vent timeout, quorum
  verdict\\
\item
  \textbf{Legal review}: fork reviewers, weighted fold by seniority\\
\item
  \textbf{Deployment control plane}: fork environment probes and publish
  candidates, race target-resolution plans, fold to a fail-closed
  publish decision with explicit smoke-gate and host-capability
  constraints
\end{enumerate}

The recurrence is framed here as discovered rather than imposed, similar
to how \emph{Physarum} discovers high-fit transport networks without
centralized planning.

In open-source/aeon-forge, this deploy-control-plane surface is
exercised by executable Bun test harnesses covering remote publish
target resolution (Nx-first, Wrangler fallback only when safe),
production smoke-gate enforcement, host compatibility constraints,
substrate capability validation, AeonPID directory
registration/current-token semantics, build-timeout cleanup, watcher
retry/debounce behavior, and metric-analyzer anomaly detection

. These checks support operational correctness claims for deployment
orchestration; they are not throughput-superiority claims.

\subsubsection{\texorpdfstring{\textbf{12.3 Wire Format
Bridge}}{12.3 Wire Format Bridge}}\label{wire-format-bridge}

The engine includes a wire format bridge to the Aeon Flow protocol. The
same 10-byte frame header (§8.2) encodes
WorkFrame\textless T\textgreater{} objects for network transmission.
Frames encoded by Aeon Pipelines transcode into frames in Aeon Flow, and
vice versa. The computation topology is independent of the transport
topology.

\subsection{\texorpdfstring{\textbf{13.
Validation}}{13. Validation}}\label{validation}

The claims are backed by executable tests across five independent
suites:

For auditability, the primary evidence-bounded claims map directly to
primary harness/artifact pairs:

\textbf{Claim family} \textbar{} \textbf{Primary harness} \textbar{}
\textbf{Primary artifacts} \textbar{}

Wall-clock matrix \textbar{} scripts/gate1-wallclock-matrix.ts
\textbar{} artifacts/gate1-wallclock-*.\{json,md\} \textbar{}

Protocol corpus \textbar{} scripts/gate2-protocol-corpus.ts \textbar{}
artifacts/gate2-protocol-corpus.\{json,md\} \textbar{}

Compression corpus \textbar{} scripts/gate3-compression-corpus.ts
\textbar{} artifacts/gate3-compression-corpus.\{json,md\} \textbar{}

Out-of-sample \textbar{} scripts/gate4-rqr-holdout.ts \textbar{}
artifacts/gate4-rqr-holdout.\{json,md\} \textbar{}

Biological effect-size mapping \textbar{}
scripts/gate5-bio-effect-size.ts \textbar{}
artifacts/gate5-bio-effect-size.\{json,md\} \textbar{}

\begin{itemize}
\tightlist
\item
  \textbf{Companion obligations and executable proofs}: pipeline
  topology, queueing containment (including exhaustive finite-trace
  work-conserving discipline coverage, representative discretized
  service-time families, a mechanized TLA+ sample-path conservation
  module for the bounded single-server case, a mechanized bounded
  multi-class open-network conservation module over finite service-law
  scenarios, a mechanized finite-support stochastic-mixture queueing
  module with positive scenario masses plus weighted-expectation checks,
  mechanized exact finite-state probabilistic queue and multiclass
  open-network kernels with distribution-level conservation invariants
  plus an explicit worst-case small-data ramp-up branch, a mechanized
  larger exact finite-support three-arrival open-network cube, and Lean
  truncation-balance theorems plus constructive infinite-weighted-sum,
  countably supported stochastic PMF, measure-theoretic lintegral,
  monotone truncation-to-limit, stable M/M/1 stationary-occupancy, and
  long-run Cesaro queueing theorems, alongside a higher-level
  queue-limit schema for stronger uninstantiated support/stability
  assumptions), flow-frame invariants, compression race properties,
  shootoff reproductions, wall-clock matrix runs across loopback
  stage-server cells and external non-loopback pools (including a
  six-distinct-host matrix; p50/p95 summaries, bootstrap confidence
  intervals, and explicit verdict artifacts), seeded heterogeneous
  protocol-corpus artifacts comparing Aeon Flow vs HTTP/3 across
  predeclared environment cells with bootstrap-CI and per-site win-rate
  criteria
  (companion-tests/artifacts/gate2-protocol-corpus.\{json,md\}), seeded
  heterogeneous compression-corpus artifacts comparing topological
  per-chunk racing against fixed-codec and heuristic baselines with
  bootstrap-CI and win-rate criteria
  (companion-tests/artifacts/gate3-compression-corpus.\{json,md\}),
  out-of-sample screening artifacts with fixed train/holdout split rules
  plus predeclared CI/threshold criteria
  (companion-tests/artifacts/gate4-rqr-holdout.\{json,md\}), comparative
  biological effect-size artifacts across predeclared condition pairs
  with Monte Carlo uncertainty propagation plus pooled bootstrap-CI
  criteria
  (companion-tests/artifacts/gate5-bio-effect-size.\{json,md\}),
  finite-DAG decomposition coverage (including edge-cover exactness and
  full source-to-sink path-set preservation), §7 formula checks
  (Worthington Whip , Speculative Tree , turbulent multiplexing
  idle-fraction bounds), quantum-topology claims (Grover-style scaling,
  Kronig-Penney band gaps as , and the linear-path-sum vs
  nonlinear-selection boundary on the path-integral correspondence,
  including same-path-family fold ablations, fixed-parameter
  toy-attention behavioral ablations with bootstrap intervals, a seeded
  Gnosis cancellation benchmark, a seeded Gnosis mini-MoE routing
  benchmark, and an artifact-generated correspondence-boundary figure),
  map/reduce readiness diagnostics (boundedness/monotonicity,
  nonzero-opportunity necessity in migration simulation, independent
  migration-simulator rank ordering, and high-readiness counterexample
  families showing non-automatic quantum asymptotics), convergence
  simulation under the three constraints, evidence-table deficits
  (including T+2 settlement under both core and broad-scope lockup
  scenarios), evidence-traceability calibration/provenance/reference
  checks, self-hosted formal artifact parsing/round-trip validation with
  aeon-logic, a parser shootoff benchmark against Java SANY
  startup-parse baselines (stabilized multi-sample harness: 9 measured
  samples after warmup, aeon-logic median 49.51 ms for 19,200 artifacts
  with IQR 48.21--49.94 ms = 387,780.9 artifacts/s; Java SANY median
  116.45 ms on BandGapVoid.tla with IQR 115.13--122.08 ms, implying
  approximately 45,156.7x normalized per-artifact throughput in this
  startup-parse harness and normalization scheme, not an end-to-end
  verification-speed claim), plus a differential parse-equivalence
  harness against SANY outcomes (100\% agreement on the current formal
  corpus for original modules, round-tripped modules and invalid-corpus
  rejections). The parser result is therefore speed plus capability
  surface: unlike the parser-only baseline, aeon-logic also exposes
  superposition chains, quorum temporal operators, topology bridges,
  Lean-sandbox project/build verification, and embedded model-checker
  interfaces in the same runtime\\
  . Mechanized TLA+ model checking across the current formal module set
  (C1--C4, queueing sample-path conservation, bounded multi-class
  queueing-network conservation, finite-support stochastic
  queueing-mixture conservation, exact finite-state probabilistic queue
  and multiclass-network kernels, larger exact finite-support
  queueing-network cubes, §7 formulas, cross-shard crossover,
  scheduler-overhead bounds, protocol/settlement deficits, quantum
  deficit identity, band-gap void, beauty-optimality scaffold), and a
  Lean 4 theorem package with constructive identities, infinite-support,
  countably supported stochastic, measure-theoretic, stable M/M/1, and
  Cesaro queueing lifts, plus explicit-assumption theorem schemas
  (including the stronger correspondence-boundary property-negation and
  general nonadditive-fold impossibility theorem plus the global
  convergence schema) verify the strongest operational claims section by
  section\\
  .\\
\item
  \textbf{Open-source flow + compression runtime}: @affectively/aeon
  flow/compression tests verify 10-byte self-describing flow frames, UDP
  fragmentation/ACK behavior, frame reassembly, flow protocol semantics,
  WASM force-mode/error semantics, and topological compression
  properties\\
  .\\
\item
  \textbf{Open-source topology engine}: @affectively/aeon-pipelines
  tests cover fork/race/fold/vent primitives, fold strategies,
  Reynolds/backpressure/turbulent multiplexing, quantum modalities,
  flow-bridge wire compatibility, domain scenarios and microbenchmarks\\
  .\\
\item
  \textbf{Open-source topology analyzer suite}:
  TopologyAnalyzer/TopologySampler tests in @affectively/aeon validate
  Betti extraction, diagnostics, void detection and executable
  protocol-topology contrasts\\
  .\\
\item
  \textbf{Open-source deployment control plane}: @affectively/aeon-forge
  Bun test suites validate remote publish planning/gating (including
  production smoke-gate constraints), host capability and substrate
  requirement checks, AeonPID directory invariants, build-timeout and
  watcher-retry behavior, and telemetry metric-analyzer anomaly
  detection\\
  . In the targeted reproducibility slice reported here (remote-publish,
  host-compat, substrate, aeonpid-directory, metric-analyzer,
  build-timeout, watcher-retry), 57 tests passed with 132 assertions.
\end{itemize}

Pass/fail totals are available from the linked suites via their
reproducible commands; parser-validated formal artifacts, mechanized
Lean theorem builds and mechanized TLC runs are all part of that
reproducible surface

.

\subsection{\texorpdfstring{\textbf{14.
Limitations}}{14. Limitations}}\label{limitations}

\textbf{Benchmark substrate.} The §7.1 step-count table remains a
topology-depth model and should not be read as a universal latency
constant. Live distributed wall-clock matrices already establish the
manuscript's deployment-facing claim at its declared benchmark scope:
loopback runs, external non-loopback single-host runs, and external
non-loopback six-host runs each satisfy predeclared improvement criteria
under impairment injection with uncertainty intervals. Protocol and
compression corpus matrices likewise establish seeded heterogeneous
corpus-scope gains under predeclared scoring rules and passing
uncertainty-interval criteria. The manuscript therefore makes
benchmark-family and seeded-corpus claims, not universal
deployment-level latency claims; additional regions, providers,
topologies, or live traffic corpora would broaden that envelope rather
than rescue the current claim.

\textbf{Independence and archival provenance.} Most artifacts in this
evidence stack are produced by self-authored open-source suites and
web-hosted companion outputs

. Independent third-party reruns, immutable archival snapshots (DOI +
content hash), and blinded cross-team replication are not yet part of
the evidence surface.

\textbf{Exemplar-selection scope.} Cross-domain biological and physical
examples are selected exemplars used to test structural correspondence
under explicit assumptions; they are not an exhaustive survey of all
candidate systems. A systematic counterexample catalog remains future
work.

\textbf{Biological effect-size substrate.} Comparative biological effect
sizes are derived from predeclared quantitative ranges already stated in
§1 (saltatory conduction velocity contrast, photosynthesis
step-vs-system efficiency contrast, and Okazaki-fragment chunk-size
contrast), with uncertainty propagation reported in
companion-tests/artifacts/gate5-bio-effect-size.\{json,md\}. This
supports bounded comparative statements for those listed conditions and
does not constitute preregistered wet-lab causal inference.

\textbf{Cross-shard cost.} The Worthington Whip crossover is
characterized in finite bounded models: after full sharding, nonzero
correction cost makes additional shards non-improving within explored
bounds (TLA+ WhipCrossover + Lean theorem + executable tests). Extending
this characterization to richer timing/service distributions and
adaptive sharding policies remains future work.

\textbf{Formal model scope.} C1--C4, queueing sample-path conservation
for finite work-conserving single-server traces, bounded multi-class
open-network queueing conservation over finite service-law scenarios,
finite-support stochastic queueing-mixture conservation in expectation,
exact finite-state probabilistic queue kernels, exact finite-state
probabilistic multiclass open-network kernels, larger exact
finite-support multiclass open-network cubes, §7 formulas (including
cross-shard crossover), scheduler-overhead bounds, protocol/settlement
deficits, quantum deficit identity, the linear-additive vs
nonlinear-selection correspondence boundary, band-gap void,
beauty-optimality scaffolds, finite-prefix truncation balance, infinite
weighted-sum queue balance, countably supported stochastic PMF queue
balance, measure-theoretic lintegral queue balance, monotone
truncation-to-limit queue balance, stable M/M/1 stationary occupancy
with finite mean, a finite-node product-form open-network occupancy law
with exact singleton mass and total mean occupancy under a supplied
stable throughput witness satisfying the traffic equations, explicit
\texttt{constructiveThroughput} and \texttt{spectralThroughput}
witness-entry paths for that Jackson layer, constructive state-dependent
queue-family wrappers for vacations, retrials, reneging, and
adaptive-routing kernels, a bounded raw two-node adaptive rerouting
witness with derived ceiling kernel, throughput bound, and linear drift
witness, an assumption-parameterized state-dependent open-network
stability/terminal-balance schema, trajectory-level Cesaro balance for
unbounded open-network sample paths, higher-level queue-limit schema,
and convergence schema are mechanized in a two-layer stack: finite-state
transition models in TLA+ (TLC), plus Lean theorems with explicit
assumptions for quantitative identities and theorem schemas for global
claims, all preflighted through the self-hosted aeon-logic parser and
Lean sandbox

. The self-verification (§10) is still scoped on the operational side to
finite, untimed state spaces with either finite-support scenario
mixtures, exact finite-state probability-mass propagation, or exact
finite-support arrival cubes; the constructive unbounded lift now covers
nonnegative measurable observables, countably supported stochastic laws,
stable M/M/1 stationarity, a traffic-equation-witness product-form layer
with explicit \texttt{constructiveThroughput} and
\texttt{spectralThroughput} witness-entry paths, constructive
state-dependent queue-family wrappers for vacations, retrials, reneging,
and adaptive routing, a bounded raw two-node adaptive rerouting witness
with derived ceiling kernel, throughput bound, and linear drift witness,
and long-run Cesaro balance under explicit convergence hypotheses rather
than arbitrary state-dependent open-network semantics; and the
correspondence-boundary proof is additionally scoped to a minimal
integer-valued fold model rather than full complex-amplitude quantum
dynamics. The remaining formal boundary is therefore narrower: richer
timing/service distributions, arbitrary exact multiclass/open networks
beyond the current bounded witnesses, automatic discharge of
traffic-equation and drift side conditions from raw network data, fully
generic positive-recurrence derivations for arbitrary unbounded
state-dependent open stochastic networks, and real-time systems (strict
latency bounds).

\textbf{Queueing theory subsumption scope.} Containment is proved for
canonical constructions (Little's Law boundary case, Erlang-style
blocking behavior and Jackson-style bottleneck limits) and extended by
executable sample-path checks that, on selected finite tick traces,
exhaustively enumerate work-conserving single-server disciplines,
representative discretized service-time families, bounded multi-class
open-network conservation over finite service-law scenarios,
finite-support stochastic arrival/service/routing mixtures in
expectation, an exact finite-state probabilistic transition kernel for a
bounded single-server queue, an exact finite-state probabilistic
multiclass open-network kernel whose worst branch already exhibits the
small-data ramp-up pathology, and a larger exact three-arrival
three-class three-node witness over the full 64-branch arrival cube

. On the unbounded side, the companion now includes constructive
finite-prefix balance theorems, infinite weighted-sum queue balance,
direct countably supported stochastic queue laws via PMF,
measure-theoretic lintegral conservation, monotone truncation-to-limit
theorems, the stable M/M/1 geometric stationary law with finite mean
queue length, a finite-node product-form open-network occupancy law with
exact singleton mass and total mean occupancy under a supplied stable
throughput witness satisfying the traffic equations, explicit
least-fixed-point and spectral witness-entry paths for that Jackson
layer, constructive state-dependent queue-family wrappers for vacations,
retrials, reneging, and adaptive-routing kernels, a bounded raw two-node
adaptive rerouting witness with derived ceiling kernel, throughput
bound, and linear drift witness, an assumption-parameterized
Foster-Lyapunov/irreducibility schema for state-dependent stationary and
terminal queue balance, and a long-run Cesaro balance theorem for
unbounded open-network sample paths under vanishing residual open age.
The remaining frontier is automatic side-condition discharge from raw
network data, arbitrary exact probabilistic multiclass/open networks
beyond the current bounded witnesses, and fully generic constructive
traffic-equation or positive-recurrence proofs for unbounded open
stochastic networks.

\subsubsection{\texorpdfstring{\textbf{14.1 Evidence-Bounded
Claims}}{14.1 Evidence-Bounded Claims}}\label{evidence-bounded-claims}

Each strong claim in this manuscript is stated with an explicit evidence
boundary and reproducible artifact path.

\begin{enumerate}
\def\labelenumi{\arabic{enumi}.}
\tightlist
\item
  \textbf{Broad deployment wall-clock claim (fixture scope): supported
  for the benchmark family.} Loopback and predeclared external
  non-loopback matrices, including a six-distinct-host run
  (workers-dev-external-multihost6-distinct), satisfy predeclared
  criteria with p50/p95 and bootstrap confidence intervals in
  gate1-wallclock-matrix.\{json,md\} and
  gate1-wallclock-external-multihost.\{json,md\}. In the six-host
  external matrix, 8/8 primary cells satisfy the primary criteria;
  across all cells, median speedup ranges 11.785x-21.620x, and the
  minimum 95\% CI lower bounds remain positive (11.365x speedup and
  3,560.98 ms latency improvement). This claim is bounded to this
  benchmark family and is not a universal production-network claim.\\
\item
  \textbf{Protocol corpus advantage claim (simulated corpus scope):
  supported for the seeded corpus family.}
  companion-tests/artifacts/gate2-protocol-corpus.\{json,md\} reports
  144 sites and 12,371 resources with 6/6 primary environment cells
  satisfying predeclared criteria; framing median gain is 72.252\% (CI
  low approximately 72.19\%); primary-cell completion-median CI lows are
  20.24-83.38 ms and completion-p95 CI lows are 19.99-98.22 ms; per-site
  win rates are 100\% on all three metrics. This claim is bounded to the
  seeded simulation corpus and does not assert internet-wide superiority
  on live traffic.\\
\item
  \textbf{Compression corpus advantage claim (seeded corpus scope):
  supported for the seeded corpus family.}
  companion-tests/artifacts/gate3-compression-corpus.\{json,md\} reports
  90 samples and 20,133,761 bytes with 4/4 primary family cells
  satisfying predeclared criteria. Primary-cell median gain vs best
  fixed codec is positive with positive CI lows (approximately
  0.0009\%-0.0075\%); median gain vs heuristic baseline is
  0.777\%-46.366\% with CI lows approximately 0.386\%-39.449\%;
  per-sample win rates are 100\% against both comparators in all primary
  cells. This claim is bounded to the seeded corpus family and does not
  assert universal superiority on live production payloads.\\
\item
  \textbf{Out-of-sample} \textbf{predictive-screening claim (model
  scope): supported for the tested simulator family.}
  companion-tests/artifacts/gate4-rqr-holdout.\{json,md\} reports
  independent train/holdout validation with predeclared scoring
  criteria: Spearman CI low 0.446, slope CI low 0.427, quartile-delta CI
  low 0.100, predictor-correlation CI low 0.176, and decile monotonicity
  violations 1 \textless= 3. This claim is bounded to the tested
  simulator family and does not assert real-world deployment
  predictivity.\\
\item
  \textbf{Biological effect-size mapping claim (predeclared
  range-extraction scope): supported as internal consistency evidence
  for the listed comparative set.}
  companion-tests/artifacts/gate5-bio-effect-size.\{json,md\} reports
  three primary biological condition pairs with positive
  uncertainty-bounded effect sizes (minimum primary-pair ratio CI low
  5.829x; median pair ratio 21.524x; pooled log-ratio 3.280 with 95\% CI
  2.289-4.360). This claim is bounded to those predeclared
  manuscript-range pairs and does not assert independent dataset
  validation or preregistered cross-lab causal inference.
\end{enumerate}

\subsection{\texorpdfstring{\textbf{15.
Conclusion}}{15. Conclusion}}\label{conclusion}

I began with a child handing a ball to another child in a line. Four
hundred handoffs. I end with a claim about structure: \textbf{failure
itself has geometry}.

The path between those two points is fork/race/fold: four operations
that express the finite DAG classes modeled in this paper.

\begin{enumerate}
\def\labelenumi{\arabic{enumi}.}
\tightlist
\item
  \textbf{Fork} raises and injects curvature.\\
\item
  \textbf{Race} moves the system through a drift field.\\
\item
  \textbf{Fold} pays reconciliation debt to return to a stable sink.\\
\item
  \textbf{Vent} maps lost branches into a bounded Void so local failure
  does not become global collapse.
\end{enumerate}

Failure is not a binary event -- broken versus working. It is a
\textbf{topological coordinate}. A system fails when its transition
kernel no longer points inward toward the small set of safe states. In
queue-like terms, that is the moment arrival pressure exceeds
restorative gravity . In topological terms, it is the moment that
unmanaged deficit turns into outward drift.

That yields four concrete lessons.

First, \textbf{failure is conserved entropy}. A race does not ``resolve
for free.'' In any finite single-survivor collapse, the Void must record
a payment. \texttt{THM-FAIL-MINCOST} is the manuscript's statement of
that floor: if you begin with live branches and end with one, at least
branch outcomes must be vented, merged, deferred, or otherwise accounted
for. Reliability is therefore not the avoidance of failure. It is the
discipline of choosing how failure is paid.

Second, \textbf{failure is gravity lost}. The drift schema shows that
stability is just a potential well. When negative drift holds outside a
bounded small set, recovery is expected. When it does not, the system
escapes. This is not metaphorical hand-waving; it is a measurable
margin.

Third, \textbf{failure is a modulator, not a stop sign}. FMFV and the
reneging-style examples teach the same lesson: systems remain reliable
when failure signals are allowed to reshape flow before the entire graph
destabilizes. The reliable system is not fail-safe. It is
\textbf{fail-flow}.

Fourth, \textbf{the Void is mapped}. It is not mystery, bad luck, or
folklore. It is the accounting ledger for discarded branches, absorbed
disagreement, and irrecoverable alternatives. Once the Void is a
coordinate system, collapse stops looking accidental.

This is the manuscript's practical achievement. It moves system design
away from reactive engineering -- patching bugs after collapse -- and
toward geometric derivation -- proving that the graph returns. The Bule
() is the stress unit of that geometry. The drift schema is its
stability test. FMFV is its validation method.

The paper does not claim a full physical unification theory. It proposes
a bounded computational-shape hypothesis with explicit executable and
mechanized scope limits. Within that scope, the evidence supports a
clear shift in viewpoint: failure is not an exception to the system.
\textbf{Failure is the shape the system takes when its restorative
geometry is too weak.}

Reliability, then, is geometry. The Newtonian step is not that bugs
vanish. It is that collapse and recovery can now be treated as lawful
state transitions on a measurable manifold.

This framing is intended as an operational modeling lens for
computation, not a replacement for physical theory.

Within this paper's modeled scope, fork/race/fold + vent is sufficient.

\subsection{\texorpdfstring{\textbf{References}}{References}}\label{references}

A. Tero, S. Takagi, T. Saigusa, K. Ito, D. P. Bebber, M. D. Fricker, K.
Yumiki, R. Kobayashi, T. Nakagaki, ``Rules for Biologically Inspired
Adaptive Network Design,'' \emph{Science}, 327(5964):439--442, 2010.

T. W. Buley, ``Aeon Pipelines: A Computation Topology Engine,''
open-source implementation, 2026.
https://forkracefold.com/content/pipeline-topology.test.ts.txt

D. Akita, I. Kunita, M. D. Fricker, S. Kuroda, K. Sato, T. Nakagaki,
``Experimental Models for Murray's Law,'' \emph{Journal of Physics D:
Applied Physics}, 50(2):024001, 2016.

J. Lobry, ``Asymmetric Substitution Patterns in the Two DNA Strands of
Bacteria,'' \emph{Molecular Biology and Evolution}, 13(5):660--665,
1996.

G. S. Engel, T. R. Calhoun, E. L. Read, T.-K. Ahn, T. Mančal, Y.-C.
Cheng, R. E. Blankenship, G. R. Fleming, ``Evidence for Wavelike Energy
Transfer Through Quantum Coherence in Photosynthetic Systems,''
\emph{Nature}, 446(7137):782--786, 2007.

J. D. C. Little, ``A Proof for the Queuing Formula: ,'' \emph{Operations
Research}, 9(3):383--387, 1961.

J. R. Jackson, ``Jobshop-Like Queueing Systems,'' \emph{Management
Science}, 10(1):131--142, 1963.

T. W. Buley, ``Aeon Core Runtime (Flow + Compression) and Test Suite,''
open-source implementation, 2026. https://github.com/affectively-ai/aeon

T. W. Buley, ``Fork/Race/Fold Companion Tests,'' reproducibility suite,
2026.
https://github.com/affectively-ai/aeon/tree/main/docs/ebooks/145-log-rolling-pipelined-prefill/companion-tests

R. P. Feynman, A. R. Hibbs, ``Quantum Mechanics and Path Integrals,''
McGraw-Hill, 1965.

J. N. Bryngelson, J. D. Onuchic, N. D. Socci, P. G. Wolynes, ``Funnels,
Pathways, and the Energy Landscape of Protein Folding: A Synthesis,''
\emph{Proteins}, 21(3):167--195, 1995.

L. Lamport, \emph{Specifying Systems: The TLA+ Language and Tools for
Hardware and Software Engineers}, Addison-Wesley, 2002.

T. W. Buley, ``Aeon Logic: Fork/Race/Fold Temporal Logic Engine and
TLC/TLA Compatibility Layer,'' open-source implementation, 2026.
https://github.com/affectively-ai/aeon-logic

Lean FRO Team, ``The Lean Theorem Prover (Lean 4),'' software and
documentation, 2026. https://lean-lang.org

T. W. Buley, ``Gnosis: A Topological Programming Language with
Self-Hosting Compiler,'' open-source implementation, 2026.
https://github.com/affectively-ai/gnosis

EURORDIS-Rare Diseases Europe, ``The Diagnosis Odyssey of People Living
with a Rare Disease: Survey overview,'' Rare Barometer report, 2024.
https://www.eurordis.org/wp-content/uploads/2024/05/Diagnosis-Survey-overview-1.pdf

Depository Trust \& Clearing Corporation (DTCC), ``DTCC 2024 Annual
Report,'' 2025. (NSCC average daily transaction value: \$2.219 trillion)
https://www.dtcc.com/annuals/2024/

T. W. Buley, ``Aeon Forge: Deployment and Routing Primitives with
Bun-Tested Control-Plane Invariants,'' open-source implementation, 2026.
https://github.com/affectively-ai/aeon-forge

L. M. N. Wu, A. Williams, A. Delaney, D. L. Sherman, P. J. Brophy,
``Increasing Internodal Distance in Myelinated Nerves Accelerates Nerve
Conduction to a Flat Maximum,'' \emph{Current Biology},
22(20):1957--1961, 2012.

I. Tasaki, ``The electro-saltatory transmission of the nerve impulse and
the effect of narcosis upon the nerve fiber,'' \emph{American Journal of
Physiology}, 127(2): 211--227, 1939.

E. Voita, D. Talbot, F. Moiseev, R. Sennrich, I. Titov, ``Analyzing
Multi-Head Self-Attention: Specialized Heads Do the Heavy Lifting, the
Rest Can Be Pruned,'' \emph{Proceedings of the 57th Annual Meeting of
the Association for Computational Linguistics}, 2019.

H. Edelsbrunner, D. Letscher, A. Zomorodian, ``Topological Persistence
and Simplification,'' \emph{Discrete \& Computational Geometry},
28:511--533, 2002.

S. Mac Lane, ``Natural Associativity and Commutativity,'' \emph{Rice
University Studies}, 49(4): 28--46, 1963.

J. D. C. Little, S. C. Graves, ``Little's Law,'' in \emph{Building
Intuition: Insights From Basic Operations Management Models and
Principles}, Springer, 2008.

C. A. Petri, ``Kommunikation mit Automaten,'' doctoral dissertation,
University of Bonn, 1962.

R. Milner, \emph{Communicating and Mobile Systems: The Pi-Calculus},
Cambridge University Press, 1999.

R. M. Tomasulo, ``An Efficient Algorithm for Exploiting Multiple
Arithmetic Units,'' \emph{IBM Journal of Research and Development},
11(1):25--33, 1967.

J. E. Smith, G. S. Sohi, ``The Microarchitecture of Superscalar
Processors,'' \emph{Proceedings of the IEEE}, 83(12):1609--1624, 1995.

M. Castro, B. Liskov, ``Practical Byzantine Fault Tolerance,''
\emph{OSDI}, 1999.

M. Yin, D. Malkhi, M. K. Reiter, G. Golan-Gueta, I. Abraham, ``HotStuff:
BFT Consensus with Linearity and Responsiveness,'' \emph{PODC}, 2019.

R. E. Blankenship, D. M. Tiede, J. Barber, G. W. Brudvig, G. Fleming, M.
Ghirardi, M. Gunner, W. Junge, D. M. Kramer, A. Melis, T. A. Moore, A.
L. Moore, J. V. Moser, D. G. Nocera, A. Nozik, D. R. Ort, W. W. Parson,
R. C. Prince, R. T. Sayre, ``Comparing Photosynthetic and Photovoltaic
Efficiencies and Recognizing the Potential for Improvement,''
\emph{Science}, 332(6031):805--809, 2011.

S. Balakrishnan, R. A. Bambara, ``Okazaki Fragment Metabolism,''
\emph{Cold Spring Harbor Perspectives in Biology}, 5(2):a010173, 2013.

S. G. Waxman, ``Determinants of Conduction Velocity in Myelinated Nerve
Fibers,'' \emph{Muscle \& Nerve}, 3(2):141--150, 1980.

G. M. Amdahl, ``Validity of the Single Processor Approach to Achieving
Large-Scale Computing Capabilities,'' \emph{AFIPS Spring Joint Computer
Conference}, 30:483--485, 1967.

J. L. Gustafson, ``Reevaluating Amdahl's Law,'' \emph{Communications of
the ACM}, 31(5):532--533, 1988.

C. E. Shannon, ``A Mathematical Theory of Communication,'' \emph{Bell
System Technical Journal}, 27(3):379--423, 1948.

J. Dean, S. Ghemawat, ``MapReduce: Simplified Data Processing on Large
Clusters,'' \emph{OSDI}, 2004.

L. K. Grover, ``A Fast Quantum Mechanical Algorithm for Database
Search,'' \emph{Proceedings of the 28th Annual ACM Symposium on Theory
of Computing (STOC)}, 212--219, 1996.

P. W. Shor, ``Algorithms for Quantum Computation: Discrete Logarithms
and Factoring,'' \emph{Proceedings of the 35th Annual Symposium on
Foundations of Computer Science (FOCS)}, 124--134, 1994.

\subsection{\texorpdfstring{\textbf{Reproducibility}}{Reproducibility}}\label{reproducibility}

Source code, test suites and protocol comparison benchmarks are
available under open-source license. The scheduler, flow protocol,
compression subsystem, computation topology engine, deploy-control-plane
invariants, formal parser/tooling layer and topological programming
language are independently testabl

\end{document}
